\documentclass[twocolumn,10pt]{article}
\usepackage[margin=1.8cm,top=2.2cm,bottom=2.2cm]{geometry}
\usepackage{mathptmx}
\usepackage[scaled=0.92]{helvet}
\usepackage{graphicx}
\usepackage{amsmath,amssymb}
\usepackage{natbib}
\setcitestyle{authoryear,round,aysep={}}
\usepackage{fancyhdr}
\usepackage[colorlinks=true,linkcolor=blue,citecolor=blue,urlcolor=blue]{hyperref}
\usepackage{booktabs}
\setlength{\parskip}{0pt}
\pagestyle{fancy}
\fancyhf{}
\fancyhead[LO,RE]{\small\itshape Time-domain SGL imaging of cloudy exoplanets}
\fancyhead[RO,LE]{\small\itshape Sontakke}
\fancyfoot[C]{\small\thepage}
\providecommand{\arxiv}[1]{arXiv:#1}
\newcommand{\acknowledgments}{\section*{Acknowledgments}}
% All text added or altered in the revision is wrapped in \rev{}, which renders
% in BOLDFACE as the editor requested.  Set \revflag=0 to produce an unmarked
% (clean) version without re-editing the source.
\newif\ifrevflag \revflagfalse
% \rev marks text added or altered in this revision in BOLDFACE, as the editor
% requested.  It is deliberately a GROUP rather than \textbf so that a single
% \rev{...} can cover blank lines, itemize environments, tables and whole
% subsections.  The \ifmmode guard makes \rev a no-op inside mathematics:
% \bfseries is illegal in math mode, and a global \mathversion{bold} would
% depend on bold math fonts for every family used.  Changed equations are
% therefore marked by the bold sentence that introduces them and are listed by
% number in the point-by-point response.  \revflagfalse yields the clean version.
\ifrevflag
  \newcommand{\rev}[1]{\relax\ifmmode{#1}\else{\bfseries #1}\fi}
\else
  \newcommand{\rev}[1]{#1}
\fi

\begin{document}

\twocolumn[{%
\begin{center}
{\Large\bfseries Time-Domain Imaging of Rotating, Cloud-Covered Exoplanets
with the Solar Gravitational Lens: Reconstruction Algorithms,
Observing-Cadence Requirements, and Mission Targets\par}
\vspace{2.5mm}
{\large Sanskar Sontakke\par}
\vspace{1mm}
{\small\itshape Independent Researcher\\
sanskarsontakke@gmail.com\quad ORCID: 0009-0008-4531-5707\par}
\vspace{2mm}
\vspace{2mm}
\begin{minipage}{0.84\textwidth}
\small
\textbf{ABSTRACT}\\[1mm]
The solar gravitational lens (SGL) can, in principle, deliver multipixel
images of terrestrial exoplanets from heliocentric distances
$z\simeq650$--900~AU. In practice the planet rotates and its cloud cover
evolves while a spacecraft raster-scans the kilometer-scale image cylinder
one ``pixel'' at a time, so the measurements are temporally aliased; recent
studies identify this dynamic inversion problem as a leading unsolved
challenge for the concept. We present an open, end-to-end simulation
framework that reproduces the published SGL photometric benchmarks (the
aperture-averaged $d/4\rho$ kernel; per-sample $\mathrm{SNR}_{\rm C}=43.16$
for 1800~s dwells on an exo-Earth at 30~pc from 650~AU; the
$0.891\,D/(d\sqrt{N})$ deconvolution penalty) and then treats rotation,
orbital illumination, and stochastic, advecting clouds explicitly. We
formulate image recovery as a regularized generalized least-squares
inversion of the time-tagged sample stream (time-domain inversion, TDI,
distinguished here from detector-level time-delay integration), in which
rotation and illumination are modeled directly by the forward operator and
cloud variability enters statistically. \rev{For a rotating, cloud-free planet observed under known spin geometry, TDI
recovers the surface albedo map with Pearson $r=0.9786\pm0.0002$
($\mathrm{SSIM}=0.848\pm0.001$) over 10 paired seeds. Rotational aliasing is
tractable under characterized ephemerides; it is not an information ceiling, and
the correlation lost within the supported harmonic range is regularization, not
spin.} \rev{Clouds, not coronal photon noise, dominate the error budget:
per-sample cloud-induced variance exceeds coronal shot noise by
$163$--$440$ ($12.8$--$21.0$ in standard deviation) at $t_s=1800$~s, collapsing
fidelity to $r=0.249\pm0.010$ at 71\% realized cover under single-campaign
budgets; under this terrestrial cloud regime, neither the present pipeline nor any published method achieves faithful reconstruction in single-campaign scalar photometry.}
\rev{Crucially, surface-correlated cloud structures locked to topography constitute a severe, irreducible failure mode: phase-locking between weather and geography suppresses reconstruction fidelity to $r = 0.1444 \pm 0.0124$ (a paired loss of $\Delta r = -0.1978 \pm 0.0175$, $p_t = 3.6 \times 10^{-7}$), which no revisit cadence can mitigate because revisits resample the same contaminated rotational phase.} \rev{The observing-cadence law is measured under two explicit,
mutually exclusive resource ledgers (fixed 8~h of photons per raster position,
or fixed 89.87~d of wall clock): replacing 4 long dwells with 64 short revisits
improves $r$ from $0.049$ to $0.571$ at fixed photons, monotonically and
unanimously across six paired seeds, while the two arms agree to within the
paired standard error. Slew, settling and metrology overheads (45~s per dwell)
reduce the duty cycle from 99.4\% to 90.9\% over that range (90.5\% for the
wall-clock-constrained arm).} \rev{Ablation of
the estimator itself shows that the algorithmic gain comes from eliminating the
per-dwell-slot cloud nuisance: exact anchored profiling of the slot means lifts
$r$ by $+0.054\pm0.006$ over a fit that ignores them ($p_t=1.1\times10^{-5}$,
10/10 seeds), whereas temporal whitening contributes nothing measurable at the
fiducial cadence ($-0.003\pm0.002$, $p_t=0.13$).} \rev{The spin-ephemeris
requirement is measured rather than asserted: a $5^\circ$ pole tilt costs
$0.057$ in $r$ and a $5^\circ$ initial-phase error $0.047$, and $60^\circ$ of
phase error destroys the map entirely. We show that residual
misalignments are not refinable in flight via profile likelihood: the whitened
$\chi^2$ is exactly invariant to assumed initial phase and monotone in assumed
period over the grid tested, so the spin state is observable through image
quality but not recoverable from photometric residuals by that route.}
\rev{Surface correlation $r(\ell)$, computed with a real spherical-harmonic
transform and $\cos\vartheta$ quadrature weights, shows band means above $0.99$ up to
the $\ell\approx18$ grid ceiling of the $36\times72$ grid when cloud-free (minimum single-degree correlation $0.987$ at $\ell=17$),
and under clouds falls below $0.5$ at $\ell\simeq9$ in a five-point running mean
($\sim2200$~km for an Earth-radius planet). A land--ocean separation of
$d'=0.73\pm0.03$ (mean $\pm$ SE over 10 seeds; $\mathrm{SD}=0.081$) at $55.8\%$ realized cover exceeds every member of a 40-sample
physical cloud-only null, giving $p\le0.024$ --- the ensemble floor, not a
measurement of $0.02$.} \rev{Finally, we map requirements onto missions: we list
the five nearest confirmed potentially habitable planets around M dwarfs plus
$\tau$~Cet~e --- whose planet interpretation has since been refuted, so it is
kept here only as a nominal-proxy dynamics illustration (not an uncertainty
propagation), evaluate image-plane dynamics ($1.34$--$5.81$~km\,s$^{-1}$ of
along-track velocity change per 90~d for the five confirmed rows, $0.11$ for the
illustrative one), discuss deep-perihelion sailcraft,
planar radioisotope tiles and optical relays, and swarms as an asymptotic
architectural scenario.}
\\[2mm]
\textbf{Keywords:} gravitational lensing: strong --- planets and
satellites: terrestrial planets --- techniques: image processing ---
methods: numerical --- space vehicles: instruments
\end{minipage}
\end{center}
\vspace{6mm}
}]



\section{Introduction}
\label{sec:intro}

Resolving the surface of an Earth-like exoplanet is beyond any credible
filled-aperture or interferometric architecture: a $10\times10$ surface
map of an exo-Earth at 10~pc requires resolution elements of
$\sim$0.85~$\mu$as, i.e., a $\sim$10$^2$~km optical aperture, and
integration times for conventional concepts reach $10^4$--$10^5$~yr once
realistic backgrounds are included \citep{tur26direct}. The solar
gravitational lens (SGL) evades this scaling by using the Sun as the
primary optical element \citep{eshleman79,tt17}. Wave-optical treatments
show that a strong-interference region forms beyond
$z_0\simeq b^2/2r_g\simeq547.8$~AU (for rays of impact parameter
$b\simeq R_\odot$, with $r_g=2GM_\odot/c^2\simeq2.95$~km), with on-axis
amplification $\mu_0=4\pi^2 r_g/\lambda\simeq1.2\times10^{11}$ at
$\lambda=1~\mu$m and angular resolution of order $10^{-10}$~arcsec
\citep{tt17,tt19}. An Earth-radius planet at 30~pc is compressed into an
image cylinder only $D_{\rm img}\simeq1.34$~km across at $z=650$~AU; a
meter-class telescope acts as a single-pixel detector that must
physically traverse the image plane, measuring the integrated brightness
of the planetary Einstein ring at each sampling location and
reconstructing the source by deconvolution
\citep{tt20psf,toth21,tt22mnras}.

The optical side of this problem is now well characterized. The
aperture-averaged point-spread function falls off as $d/(4\rho)$, so each
sample mixes light from the entire planet; removing this blur incurs a
deconvolution penalty $\mathrm{SNR_R/SNR_C}\simeq0.891\,D/(d\sqrt{N})$
that is strongly mitigated by wide pixel spacing \citep{toth21,tt22mnras}.
The solar corona is the dominant photon background; for an exo-Earth at
$z_0=30$~pc observed from $z=650$~AU with a $d=1$~m telescope at
$\lambda=1~\mu$m, the signal and corona rates are
$Q_{\rm exo}=8.01\times10^4$~s$^{-1}$ and
$Q_{\rm cor}=6.20\times10^9$~s$^{-1}$, giving a convolved-image
signal-to-noise ratio $\mathrm{SNR_C}=43.2$ per 1800~s dwell
\citep{tt22mnras,tur26bench}. Megapixel-class static imaging within
mission-compatible integration times follows \citep{toth21,tt22mnras},
and independent analyses of ring-based approaches
\citep{mm22} sharpened the role of coronal statistics in these budgets.

The unsolved problem is temporal. The planet rotates and orbits, its
illumination changes, and---most seriously---its cloud cover reorganizes
on day timescales while the raster is acquired over weeks to months.
\citet{toth23} performed the first dynamic reconstructions of a rotating,
orbiting planet and recovered surface albedo only at
$\mathrm{SNR}\lesssim2$. \citet{toth25clouds} added a stochastic cloud
model, treated the clouds as noise to be averaged out (a
``long-exposure'' philosophy), and found that reconstructed SNR falls to $1.13$ at $6.8\%$ and $0.80$ at $13.7\%$ cloud cover (falling below unity between $7\%$ and $14\%$), concluding with ``a cautious
assessment of the practical feasibility'' of SGL imaging given that
Earth-like planets have $\sim$70\% cloud cover. The comprehensive
benchmark of \citet{tur26bench} reaches a complementary conclusion: a
static inverse applied to rotating cloudy data is the wrong estimator;
phase-registered robust coadds recover persistent surface structure but
at the cost of $\gtrsim$15 registered visits per rotational phase bin,
and a ``full spin-resolved dynamic retrieval pipeline'' is explicitly
deferred as a next-stage deliverable. The covariance-aware biosignature
framework of \citet{cov26} likewise argues that SGL inference must
propagate structured error terms rather than treat all residuals as
white. In short: the field agrees that the time-domain inverse problem is
the gate through which SGL imaging must pass, and no published pipeline
yet passes it under Earth-like cloud cover. We emphasize that the present pipeline also does not achieve high fidelity at Earth-like cloud cover ($r = 0.249 \pm 0.010$ at $71\%$ cover under single-campaign budgets in our adopted scalar-photometry raster architecture), and that no matched head-to-head comparison against the 15-visit phase-registered architecture of \citet{tur26bench} is established here because our sampling budgets explore visit counts under constrained photon and wall-clock allocations rather than the dedicated phase-binning regime. Our results are strictly scoped to the single-pixel photometric raster-scan model.

This paper attacks that gate directly, with five specific contributions.
First (Sections~\ref{sec:model}--\ref{sec:methods}), we build a transparent,
end-to-end simulator---partially validated against the published photometric
and deconvolution benchmarks---in which planetary rotation, orbital
illumination, and stochastic advecting clouds are modeled explicitly.
\rev{To that first item we add two rigorous audits: a
convergence study of the intra-dwell exposure integral (Section~\ref{sec:quad}),
which shows that the 3-node rule previously used aliases rotational harmonic
$m=36$ at unit gain, and a numerical-precision and conditioning audit of the
normal equations (Sections~\ref{sec:infaudit}, \ref{sec:precisionnum}), which
identifies the constant-albedo mode as the quantity the slot nuisance actually
regularizes.}
\rev{Second (Section~\ref{sec:results}), we formulate reconstruction as
regularized generalized least squares on time-tagged samples (time-domain
inversion, TDI; distinct from hardware time-delay integration in detector
readouts), with the per-dwell-slot cloud mean handled by exact anchored nuisance
profiling rather than approximate deflation heuristics. We isolate the
observing-cadence law under two named, mutually exclusive resource ledgers
(fixed photons per raster position, or fixed wall clock) rather than a single
unconstrained budget, and ablate the estimator --- slot handling, whitening,
anchoring --- holding data and operator fixed. Every number is reported as a
mean and standard error over paired seeds, with paired $t$-statistics, win
counts and the exact attainable significance floor, and with a physical
cloud-only null in place of geometry-scrambled surrogates.}
\rev{Third, we broaden the robustness tests to non-stationary latitudinal cloud
banding, multi-timescale weather, per-pass independent weather, surface-coupled
clouds, and a full $5\times5$ grid of affine cloud-climatology errors. Crucially, we identify surface--cloud coupling as a major failure mode ($\Delta r = -0.198$, reducing $r$ to $0.144$) that cannot be bought off by cadence, and we
state explicitly which metrics are blind to which errors ($r$ is exactly affine
invariant; bias, RMSE and contrast are not).}
\rev{Fourth, we extend the geometric requirement to spin-pole tilt and initial
phase with measured tolerances, and we demonstrate that
the residual spin state is not refinable in flight from the profile likelihood:
the whitened objective is exactly invariant to assumed initial phase.}
\rev{Fifth (Section~\ref{sec:mission}), we evaluate effective spatial resolution
$r(\ell)$ with a real spherical-harmonic transform and an explicit quadrature
error bound, land--ocean detection against two nulls with stated ensemble
floors, and translate cadence requirements into mission terms using an
effective-independent-look count that depends on the assumed weather
decorrelation time; the target table is a reproducible nominal-proxy list, not
an uncertainty propagation.}

\section{SGL Photometric Model and Validation}
\label{sec:model}

\subsection{Optical response}
We adopt the standard monopole, aperture-averaged description of SGL
image formation \citep{tt20psf,toth21,tur26bench}. A source-plane
coordinate $\boldsymbol{\xi}$ maps to the image-plane coordinate
$\boldsymbol{\rho}\simeq-(z/z_0)\,\boldsymbol{\xi}$; an Earth-radius
planet at $z_0=30$~pc observed from $z=650$~AU projects to
\begin{equation}
D_{\rm img}=2R_p\,\frac{z}{z_0}=1.338\;{\rm km}.
\end{equation}
A telescope of aperture $d$ used as a ring photometer averages the
$J_0^2$ oscillations of the monopole PSF; the effective kernel is
\begin{equation}
K(0)=1,\qquad K(\rho>0)\simeq\frac{d}{4\rho},
\label{eq:kernel}
\end{equation}
renormalized so that the mean convolved signal over the projected disk
of a reference (albedo 0.3, fully illuminated) planet is unity
\citep{tur26bench}. Our discrete kernel follows the $d/4\rho$ tail over
the full grid (Fig.~\ref{fig:model}a). Because $K$ falls only as
$1/\rho$, every sample mixes light from the entire disk---the property
that couples the temporal variability of any part of the planet into
every measurement.

\rev{\emph{Central-cell discretisation and normalization conventions.} Two numerical conventions govern Eq.~\ref{eq:kernel}, and we audit both directly from the pipeline archives. First, regarding the central singularity of the $1/\rho$ kernel: on a discrete cell of pitch $\Delta_{\rm img}=20.913$~m with aperture $d=1$~m, the exact cell-mean integral is $K_{00,\rm cellmean} = (d/\Delta_{\rm img})\ln(1+\sqrt{2}) \simeq 0.042144$ (and $0.010536$, $0.021072$, $0.084288$ at $n=16,32,128$, with the point-to-cellmean ratio reaching up to $94.9$; archive \texttt{kernel\_convergence.json}). Setting $K(0)=1.0$ evaluates the kernel at the sub-pixel point rather than its cell-averaged mean. To quantify the sensitivity of the inversion to this choice, we execute a controlled core-ablation experiment (archive \texttt{core\_ablation.json}): at $f_c=0$, the point kernel achieves $r=0.98034$ and $\mathrm{SSIM}=0.85126$, whereas the cell-mean kernel yields $r=0.75728$ and $\mathrm{SSIM}=0.48556$ --- a 22\% decrease in correlation and 43\% in SSIM from the central-cell convention alone. We adopt $K(0)=1.0$ as the intended benchmark to maintain exact alignment with the published baseline of \citet{tur26bench}, but note that central-cell discretisation is an important modeling sensitivity.}

\subsection{Photon rates and noise}
We normalize the scene such that a value of 1 corresponds to the
disk-mean signal of the reference planet, which maps to the fiducial
rates of \citet{tt22mnras}:
$Q_{\rm exo}=8.01\times10^4\,(d/1\,{\rm m})^2(650\,{\rm AU}/z)^{1/2}
(30\,{\rm pc}/z_0)(\lambda/1\,\mu{\rm m})$~s$^{-1}$ and
$Q_{\rm cor}=6.20\times10^9\,(d/1\,{\rm m})^2(650\,{\rm AU}/z)^2
(\lambda/1\,\mu{\rm m})$~s$^{-1}$ after coronagraphic rejection and
annular extraction, with the mean corona assumed subtracted and only its
shot noise retained. Each dwell of duration $t_s$ at image-plane position
$\boldsymbol{\rho}_k$ yields a Gaussian sample with standard deviation
\begin{equation}
\sigma_k=\frac{\sqrt{(Q_{\rm cor}+Q_{\rm exo}\bar{s}_k)\,t_s}}
{Q_{\rm exo}t_s},
\label{eq:noise}
\end{equation}
where $\bar s_k$ is the convolved scene value. For $t_s=1800$~s this
gives $\mathrm{SNR_C}=1/\sigma=43.16$, matching \citet{tur26bench}
exactly within the adopted photon-noise assumptions.
\rev{We clarify the normalisation convention: under Lambertian sphere geometry, the disk average of the projected planetary flux is $\langle \mu \rangle = 8/(3\pi) \simeq 0.848826$ of the central sub-stellar peak. If the signal is normalized to the central sub-stellar peak, $\mathrm{SNR_C}=43.16$; if normalized strictly to the disk-mean flux without that peak boost, the resulting convolved SNR is $\mathrm{SNR}_{C,\mathrm{no}\text{-}\mu} = 43.16 \times 8/(3\pi) = 36.6344$ (archive \texttt{validation.npz}). We adopt the sub-stellar 1.0 normalization to match the published 43.16 benchmark.}

\subsection{Deconvolution-penalty validation and numerical limits}
As a check of the discrete operator we reproduce the deconvolution
noise-amplification analysis of \citet{toth21}. For a static, fully
illuminated planet we convolve, add white noise, deconvolve by Fourier
quotient, and measure the ratio $\mathrm{SNR_R/SNR_C}$ (the ratio of
reconstructed signal-to-noise ratio to convolved input signal-to-noise ratio,
evaluated on the projected disk support, i.e., the reciprocal of the noise amplification factor $\sigma_R/\sigma_C$).
Figure~\ref{fig:validation} compares measurements with the theoretical formula
$\min(1.0, 0.891\,D/(d\sqrt{N}))$.
The theoretical curve is capped at $1.0$ (in \texttt{src/run\_experiments.py})
because unregularized deconvolution cannot increase SNR beyond the input.
At $n=32$ ($D=41.8$~m), the uncapped analytic formula yields $4\times0.2912=1.1646$,
and we measure $0.657$ vs.\ the capped $1.0$.
For $n=64$ ($D=20.9$~m), measured is $0.3244$ vs.\ analytic $0.2912$ (11\% above);
for $n=128$ ($D=10.5$~m), measured is $0.1057$ vs.\ analytic $0.0728$ (45\% above).

\rev{The sign reversal between coarse and fine grids reflects competing numerical effects.
At $n=32$, discrete pixel averaging across the central singularity acts as an effective
low-pass filter, and zero-padded circular boundary conditions attenuate noise amplification below the continuum integral (measured $<$ analytic).
Conversely, at fine pitches ($n=64$ and $128$), discrete high-wavenumber noise amplification outpaces the continuous circular integral (measured $>$ analytic) because the discrete Fourier quotient on a finite, zero-padded grid suffers from spectral leakage and ill-conditioned near-zero transfer function division near the grid Nyquist boundary, where regularisation is absent.}
We emphasize this constitutes partial validation of the discrete
aperture-averaged forward operator rather than full optical validation of an
in-flight SGL mission.

\begin{figure*}[t]
\centering
\includegraphics[width=0.92\textwidth]{figures/fig_model.pdf}
\caption{Model overview. (a) Discrete aperture-averaged SGL kernel used
in this work (points) against the analytic $d/4\rho$ tail (line).
(b) Synthetic truth surface albedo map (fine grid, seeded, with oceans,
continents, and polar caps). (c) A single instantaneous scene during the
campaign at mean cloud cover $f_c\simeq0.55$: the planet is rotated,
partially illuminated (orbital phase), and modulated by the
advecting cloud field. (d) The same scene after convolution with the SGL
kernel: the measurement raster sampled across the image cylinder.}
\label{fig:model}
\end{figure*}

\begin{figure}[t]
\centering
\includegraphics[width=0.98\columnwidth]{figures/fig_validation.pdf}
\caption{Validation of the discrete operator: measured deconvolution
noise penalty vs.\ the analytic estimate of \citet{toth21}. The $\sim$45\%
difference at $n=128$ reflects discrete pixel-averaging of the central
singularity and finite-domain boundary attenuation.}
\label{fig:validation}
\end{figure}

\section{Dynamic Scene and Campaign Model}
\label{sec:scene}

\subsection{Rotating, illuminated planet}
The truth surface is a seeded synthetic albedo map (ocean 0.06,
continents $0.32\pm0.10$, polar caps up to 0.60; 30\% land fraction)
generated by thresholded multi-octave Gaussian random fields on a
$144\times72$ longitude--latitude grid (Fig.~\ref{fig:model}b). The planet
rotates with period $P_{\rm rot}=24$~h and is illuminated by its host star
with a Lambertian terminator whose sub-stellar point drifts through
\rev{$86.1^\circ$ of orbital phase over the 87.33-day fiducial campaign
($P_{\rm orb}=1$~yr, $360^\circ\times87.33/365.25$; $88.7^\circ$ over a
90.0-day window)}.
\rev{The simulation holds the sub-stellar latitude fixed at the equator:
obliquity is not modeled, so the illumination drift here is purely orbital in
longitude and the seasonally-migrating intertropical zone of a tilted planet is
outside the tested regime (Section~\ref{sec:discussion}).}
\rev{For the short-period habitable-zone planets around M dwarfs listed in
Section~\ref{sec:targets} ($P\sim10$--20~d), a 90-day campaign spans 5--9 full
orbits and therefore samples all orbital phases repeatedly; the single
sub-90$^\circ$ drift modeled here is the \emph{least} phase-covering case, not
a conservative envelope for every target, because the two regimes differ in
which longitudes are ever illuminated (a synchronous, zero-obliquity planet
keeps a fixed sub-stellar longitude). We therefore report the tested case and
do not claim the benchmark is conservative for all listed targets.}
\rev{Each dwell is integrated over $n_{\rm sub}=16$ Gauss--Legendre
sub-exposures spanning the dwell, so intra-dwell rotational smear
($7.5^\circ$ per 1800~s dwell) is present in both the synthetic data and the
forward operator. A uniformly-spaced 3-node rule was tested and rejected: its
node spacing of $t_s/3$ aliases azimuthal harmonic $m=36$ onto the disk-integrated
mean with unit gain instead of annihilating it, a response error of order unity
that propagates directly into the reconstruction (Section~\ref{sec:quad}).}

\subsection{Stochastic cloud fields: stationary and structured}
The baseline cloud model is an advecting Ornstein--Uhlenbeck (OU) Gaussian
random field on the fine grid: spatial correlation length $12^\circ$,
zonal advection $6^\circ$\,day$^{-1}$, temporal decorrelation $\tau_c=4$~d,
mapped through a soft threshold to an opacity screen $c\in[0,1]$. The
effective scene is $A_{\rm eff}=(1-c)\,A_{\rm surf}+c\,A_{\rm cl}$ with
$A_{\rm cl}=0.65$. \rev{The sweep is over the \emph{specified} cover parameter
$f_c\in\{0,0.25,0.40,0.55,0.70\}$; because the soft threshold saturates, the
realized mean covers are slightly different ($0.237$, $0.398$, $0.558$, $0.710$
respectively, averaged over seeds) and every table in this paper quotes
specified $f_c$ together with the realized cover it produced.}

\rev{To test robustness against non-terrestrial weather we implement four
extended climatology variants, each re-normalized so that the realized mean
cover matches the baseline $f_c$ it is compared with (otherwise a variant that
simply has more or less cloud would be confounded with a variant that has a
different cloud \emph{pattern}):}
(i) \rev{{\it Banded zonal profile:} the opacity field is multiplied by
$w_{\rm lat}(\theta)=1+0.35\cos2\theta-0.25\cos4\theta$, normalized to unit
mean. We previously labelled this an ``ITCZ'' profile, which it is not: its
maximum ($1.31$) lies at $\pm35^\circ$ --- the subtropical region an ITCZ model
should \emph{suppress} --- its equatorial value is a mild $1.10$, and it is
symmetric about the equator and never migrates. It is a test of latitudinal
inhomogeneity only, and we report it as such.}
(ii) {\it Surface-correlated cloudiness:} coupling cloud opacity to surface
reservoirs, $g_{\rm eff}=g-\beta_{\rm surf}(A_{\rm surf}-\bar{A})$, enhancing
convection over low-albedo tropical oceans; and
(iii) {\it Two-timescale temporal evolution:} \rev{summing a rapid convective
mode and a persistent synoptic regime with weights $(0.3,0.7)$ and correlation
times $(1~{\rm d},10~{\rm d})$, normalized to unit variance.}
\rev{(iv) {\it Per-pass independent weather:} a fresh, temporally uncorrelated
cloud realization for every pass, included as the causal control of
Section~\ref{sec:cadence} rather than as a physical climate.}

\subsection{Campaign and sampling geometry}
The campaign matches current swarm concepts \citep{tt22mnras,helvajian23}:
$N_{\rm sc}=16$ spacecraft sample an $n\times n=64\times64$ raster
($\Delta_{\rm img}=20.9$~m pitch, 4096 pixels) in boustrophedon order,
divided into 256 dwell slots. In each dwell slot, 16 pixels are sampled
simultaneously by spacecraft spaced across adjacent raster tracks with
transverse separations of $\sim$80--300~m. \rev{Nominal dwells are
$t_s=1800$~s with $M_p=16$ complete passes, giving 65{,}536 time-tagged
samples, $8.00$~h of integration per raster position, $85.33$~d of exposure
time, $2.00$~d of charging overhead inside the campaign window, and a
wall-clock length of $\mathbf{87.33}$~d --- not the $90.0$~d quoted previously,
which silently assumed zero overhead. The overhead ledger charges $45$~s to
each of the $M_p\times256$ dwells, $2.13$~d in all, but the final pass's
overhead falls after the last sample and is not counted inside the campaign
window, so the wall clock is $85.33+2.00=87.33$~d, not $85.33+2.13$~d.
Inter-pass dead time is now an explicit,
auditable ledger instead of a random draw: the fiducial campaign charges
$T_{\rm oh}=45$~s per dwell (30~s translation slew, 10~s pointing settling,
5~s inter-spacecraft metrology) and no additional random gap.

\emph{Raster kinematics and continuous-scan vs.\ stop-and-dwell.} We explicitly distinguish stop-and-dwell stepping from continuous-scan rastering. If each step of pitch $\Delta_{\rm img} = 20.913$~m is executed as a rest-to-rest stop-and-dwell translation within the nominal 30~s translation budget, kinematics imposes a strict impulse lower bound: achieving displacement $\Delta$ in time $\Delta t$ with a symmetric bang-bang thrust profile requires peak velocity $v_{\rm max} = 2\Delta/\Delta t \simeq 1.394$~m\,s$^{-1}$ and a total velocity increment of $\Delta v_{\rm hop} = 2 v_{\rm max} \simeq 2.788$~m\,s$^{-1}$ per hop. Over $16{,}320$ neighbor transitions during the 16 passes, the cumulative impulsive $\Delta v$ would reach $\approx 22.75$~km\,s$^{-1}$ per spacecraft before inter-pass repositions, an expenditure infeasible for sub-kilogram sailcraft or small exploration probes without continuous low-thrust propulsion. Accordingly, we treat these raster positions as idealized prescribed samples along the trajectory. Operationally, SGL formations would execute continuous-scan raster tracks or smooth drifting orbits through the image cylinder with continuous exposure integration, where transverse actuation is minimized to gentle station-keeping rather than discrete stop-and-dwell maneuvers. The discrete dwell times and overheads charged here serve as an explicit photometric and temporal resource ledger, not a verified closed-loop guidance and control trajectory.}

\rev{\subsection{Resource conventions: two mutually exclusive ledgers}
\label{sec:ledger}
The referee correctly observed that ``fixed 8~h per pixel'' and ``fixed 90-day
campaign'' cannot both be satisfied as $M_p$ grows. We therefore run the
cadence study under two explicitly named arms and never mix them.
\begin{itemize}
\item {\bf Arm A (photons):} hold $M_p t_s=8$~h per raster position and let the
wall clock be whatever it needs. Then $t_s=7200,3600,1800,900,450$~s for
$M_p=4,8,16,32,64$ give $85.73,86.27,87.33,89.47,\mathbf{93.73}$~d; the
$M_p=64$ configuration \emph{violates} a 90-day cap.
\item {\bf Arm B (wallclock):} force the campaign to $89.87$~d including
overheads and let $t_s$ fall to $7548.8,3751.9,1853.4,904.2,\mathbf{429.6}$~s,
so the photons per position fall from $8.39$~h to $7.64$~h instead of staying
at 8~h.
\end{itemize}
Every cadence number in Section~\ref{sec:cadence} is reported per arm, and the
two arms answer different questions: Arm A isolates the statistical effect of
revisit count at constant photons, Arm B is what a 90-day mission can actually
buy.}

\rev{Cloud-induced signal fluctuations are calibrated against the data the
estimator actually sees, i.e.\ the per-sample standard deviation of the
cloud-only component of the disk-integrated signal with the 16-node
Gauss--Legendre exposure integration of Section~\ref{sec:quad}. The resulting
values are $\sigma_{\rm cl}=0.296,0.378,0.441,0.486$ for $f_c=0.25,0.40,0.55,
0.70$, against the coronal single-dwell noise $\sigma_{\rm cor}=0.0232$ at
$t_s=1800$~s. That is a standard-deviation ratio of $\mathbf{12.8}$--
$\mathbf{21.0}$, or a \emph{variance} ratio of $\mathbf{163}$--
$\mathbf{440}$. The previously quoted ``14--23 times'' was a standard-deviation
ratio computed with the rejected 3-node quadrature; it is neither the variance
ratio its wording implied nor reproducible under the corrected operator.
Planetary weather dominates the single-sample photometric error budget by more
than two orders of magnitude in variance.}

\section{Reconstruction Methods}
\label{sec:methods}

\subsection{Time-domain inversion (TDI)}
\label{sec:tdi}
\rev{Two notational points precede the estimator. We write the target distance
as $z_0$ (pc-scale) and the spacecraft heliocentric distance as $z$ (AU-scale),
reserving $z/z_0$ for the imaging magnification; the minimum heliocentric
distance for a ray of impact parameter $b$ is $z_{\min}\simeq b^2/2r_g$ and is
\emph{not} the same quantity as the target distance $z_0$. Second, all
reconstructions use $n_{\rm sub}=16$ Gauss--Legendre sub-exposures (Eq.~\ref{eq:forward}); the number of nodes is a parameter of the operator, not a
free-tuning knob, and Section~\ref{sec:quad} shows why the previously used
3-node rule is inadequate.}
Let $\mathbf{s}\in\mathbb{R}^{n_s}$ be the static surface albedo map on
a $72\times36$ reconstruction grid ($n_s=2592$, coarser than the truth grid).
Each sample $k$, taken at time $t_k$ and image-plane pixel $p_k$, satisfies
\begin{equation}
y_k=\mathbf{f}_k^{\sf T}\mathbf{s}+\varepsilon_k,\qquad
\mathbf{f}_k^{\sf T}=\sum_{m=1}^{n_{\rm sub}}w_m\,\mathbf{k}_{p_k}^{\sf T}\,\Pi(t_k+\delta t_m),
\label{eq:forward}
\end{equation}
where $\Pi(t)$ renders the map at the rotation phase, illumination, and
spin-pole geometry of time $t$, and $\mathbf{k}_{p_k}$ is the SGL kernel row.
\rev{$\{\delta t_m,w_m\}$ are the Gauss--Legendre nodes and weights of the dwell
interval $[-t_s/2,t_s/2]$, with the weights normalized to sum to unity, so
$\mathbf{f}_k^{\sf T}$ is the exposure-averaged kernel row rather than a
pointwise one; the equal-weight average written in v2 is not the quadrature the
data are generated with.}
Setting $n_{\rm sub}=1$ corresponds to an instantaneous midpoint model;
\rev{$n_{\rm sub}=16$ uses the Gauss--Legendre nodes $\delta t_m$ and weights of
the dwell interval, which is what the synthetic data are integrated with.}
The error $\varepsilon_k$ contains coronal shot noise and cloud fluctuations.

\rev{\emph{Estimand definition and affine debiasing.} We explicitly separate the apparent (scene-integrated) map $\mathbf{m} = (1 - f_c)\mathbf{s} + f_c A_{\rm cl}\mathbf{1}$ from the bare physical surface $\mathbf{s}$. The time-domain inverse problem solves for the effective scene map under the linear operator $\mathsf{F}$. To extract the bare physical surface $\mathbf{s}$, we apply an affine post-processing correction:
\begin{equation}
\hat{\mathbf{s}}=\frac{\hat{\mathbf{m}}-f_cA_{\rm cl}\,\mathbf{1}}{1-f_c}\,.
\label{eq:debias}
\end{equation}
Crucially, this debiasing is evaluated using the \emph{specified} nominal cover parameter $f_c$ (e.g., $0.55$ or $0.70$) and nominal cloud albedo $A_{\rm cl}=0.65$. Because thresholded Gaussian random fields produce slightly offset \emph{realized} mean covers ($0.558$ and $0.710$), an operational cover mismatch of $\Delta f_c \approx 0.008$--$0.010$ is naturally present. This minor mis-registration induces small residual albedo biases in the retrieved maps (on the order of $10^{-3}$ to $10^{-2}$), including the cloud-free baseline bias of $-0.0006\pm0.0006$ (Table~\ref{tab:clouds}).}

TDI solves the regularized generalized least-squares problem:
\begin{equation}
\hat{\mathbf{m}}=\arg\min_{\mathbf{m}}\;
\|\mathsf{C}^{-1/2}(\mathbf{y}-\mathsf{F}\mathbf{m})\|^2
+\lambda_{\rm eff}\,\|\mathsf{L}\mathbf{m}\|^2,
\label{eq:tdi}
\end{equation}
with spherical-grid Laplacian $\mathsf{L}$ and
$\lambda_{\rm eff}=\lambda\,{\rm tr}(\mathsf{F}^{\sf T}\mathsf{C}^{-1}\mathsf{F})/{\rm tr}(\mathsf{L}^{\sf T}\mathsf{L})$
($\lambda=3\times10^{-3}$, \rev{carried unchanged: it is a hardcoded
constant of the released pipeline, not a quantity re-tuned in this revision, and
no calibration seed is claimed for it --- see Section~\ref{sec:lamsens}}).
\rev{That weight is used unchanged throughout the suite; it is not
claimed to be optimal --- Section~\ref{sec:lamsens} sweeps six decades of it and
finds the objective unable to select one without truth.}
\rev{$\mathsf{C}$ is the working covariance, block diagonal over image-plane
pixels: the $M_p$ samples of pixel $p$ (one per pass, at the same raster
position) are correlated by the OU cloud field,}
\begin{equation}
\rev{(\mathsf{C}_p)_{ij}=\sigma_{\rm cl}^2\exp(-|t_i-t_j|/\tau_c)+\delta_{ij}\sigma_i^2 ,}
\label{eq:cp}
\end{equation}
\rev{while different pixels are treated as conditionally independent given the
apparent map (the physical spatial correlation length of clouds, $\sim12^\circ\approx1334$~km, is absorbed into an effective variance $\sigma_{\rm cl}^2$ rather than modeled across spacecraft tracks; Section~\ref{sec:discussion} delivers the detailed analysis of what this approximation costs).}
\rev{Four weighting/projection combinations are ablated in
Section~\ref{sec:ablation}: (White) diagonal
$\mathsf{C}=\mathrm{diag}(\sigma_k^2+\sigma_{\rm cl}^2)$, the assumption that
clouds average like uncorrelated white noise \citep{toth25clouds}; (OU)
Eq.~\ref{eq:cp} without any slot treatment; (approximate deflation), the heuristic
per-slot projection described below; and (Full TDI), exact nuisance profiling
with Eq.~\ref{eq:cp}.}

\rev{\subsection{Nuisance profiling of the simultaneous-sample offsets}
\label{sec:profiling}
Because $K(\rho)\propto1/\rho$, an instantaneous disk-integrated cloud
fluctuation perturbs all $N_{\rm sc}=16$ simultaneous samples coherently. The
v2 treatment subtracted the per-slot mean of both $\mathbf{y}$ and
$\mathsf{F}$ and rescaled the marginal noise by $\sqrt{1-1/N_{\rm sc}}$,
claiming that this ``removes 256 common modes (6.25\% rank reduction) while
retaining $1-1/N_{\rm sc}=93.75\%$ of the information''. Both the mechanism and
the numbers in that sentence are wrong, and the revised treatment replaces it
with the exact profile likelihood of the referee's eq.~(R11).}

\rev{Write the model with slot nuisance offsets $\boldsymbol\alpha$,
\begin{equation}
\mathbf{y}=\mathsf{F}\mathbf{s}+\mathsf{B}\boldsymbol\alpha+\boldsymbol\varepsilon,
\qquad \boldsymbol\varepsilon\sim\mathcal N(0,\mathsf{C}),
\label{eq:nuisance}
\end{equation}
where $\mathsf{B}$ is the indicator matrix of dwell slots ($\mathsf{B}$ has one
column per slot, 4096 columns). The key structural fact, which the v2
``orthogonal projection'' narrative missed, is that $\mathsf{C}$ is block
diagonal over \emph{image-plane pixels}, while $\mathsf{B}$ couples slots
\emph{within} a pixel. The $M_p$ samples belonging to pixel $p$ sit in $M_p$
different slots, so}
\begin{equation}
\rev{\mathsf{H}\equiv\mathsf{B}^{\sf T}\mathsf{C}^{-1}\mathsf{B}
=\mathrm{blkdiag}\left(\mathsf{H}_1,\ldots,\mathsf{H}_{n_{\rm bin}}\right),}
\label{eq:H}
\end{equation}
\rev{with $n_{\rm bin}=256$ blocks of size $M_p\times M_p$ --- \emph{not}
diagonal, i.e.\ the offsets are statistically coupled through
the shared temporal cloud covariance of a single pixel. Profiling them out
therefore requires inverting those $M_p\times M_p$ blocks: with
$\mathsf{G}_j=\mathsf{B}_j^{\sf T}\mathsf{C}^{-1}\mathsf{F}$,
$\mathbf{g}_j=\mathsf{B}_j^{\sf T}\mathsf{C}^{-1}\mathbf{y}$ and
$\mathsf{H}_j=\mathsf{B}_j^{\sf T}\mathsf{C}^{-1}\mathsf{B}_j$, the free
(unconstrained) profile is}
\begin{align}
\rev{\mathsf{A}_{\rm free}}&\rev{=\mathsf{F}^{\sf T}\mathsf{C}^{-1}\mathsf{F}
-\sum_j \mathsf{G}_j^{\sf T}\mathsf{H}_j^{-1}\mathsf{G}_j,}
\label{eq:profilea}\\
\rev{\mathbf{b}_{\rm free}}&\rev{=\mathsf{F}^{\sf T}\mathsf{C}^{-1}\mathbf{y}
-\sum_j \mathsf{G}_j^{\sf T}\mathsf{H}_j^{-1}\mathbf{g}_j\,.}
\label{eq:profileb}
\end{align}
\rev{With the standard anchoring constraint $\sum_b\alpha_b=0$ this collapses
to a rank-one correction,}
\begin{equation}
\rev{\mathsf{A}=\mathsf{A}_{\rm free}+\mathbf{q}\mathbf{q}^{\sf T}/S\,,\qquad
\mathbf{b}=\mathbf{b}_{\rm free}+\mathbf{q}\,m/S\,,}
\label{eq:anchor}
\end{equation}
\rev{where $\mathbf{v}_j=\mathsf{H}_j^{-1}\mathbf{1}$ and}
\begin{equation}
\rev{\mathbf{q}=\sum_j\mathsf{G}_j^{\sf T}\mathbf{v}_j\,,\qquad
S=\sum_j\mathbf{1}^{\sf T}\mathbf{v}_j\,,\qquad
m=\sum_j\mathbf{v}_j^{\sf T}\mathbf{g}_j\,.}
\label{eq:anchordef}
\end{equation}
\rev{which is exact and costs only the block inversions already required by
Eqs.~\ref{eq:profilea}--\ref{eq:profileb}. Without anchoring the offsets are only
projected out, and the estimator loses the disk-integrated albedo information
entirely; Eq.~\ref{eq:anchor} is what ``full TDI'' means in this paper.}

\subsection{Exposure quadrature: convergence of the smear operator}
\label{sec:quad}
\rev{Eq.~\ref{eq:forward} averages the illumination/rotation operator over the
dwell, and the synthetic data are generated with the same rule, so the
sub-exposure count $n_{\rm sub}$ is part of the operator definition rather than
a tuning knob. v2 used three \emph{equally spaced}, equal-weight sub-exposures at
$-t_s/3$, $0$, $+t_s/3$, which is neither a Gauss--Legendre rule nor a convergent
quadrature of the smear kernel. Writing the exact rotational-harmonic response of
a dwell of length $t_s$ as $R_m(t_s)=\mathrm{sinc}(m\Omega t_s/2)$ with
$\Omega=2\pi/P_{\rm rot}$, the equal-spaced rule gives}
\begin{equation}
\rev{R_{m}^{({\rm eq}\,3)}=\tfrac13\left[2\cos\!\left(\tfrac{m\Omega t_s}{3}\right)+1
\right]\,,}
\label{eq:alias}
\end{equation}
\rev{so at $t_s=7200$~s the exact response $R_{36}=0$ becomes
$R_{36}^{({\rm eq}\,3)}=1.000000$.}
\rev{At $P_{\rm rot}=24$~h the node spacing $t_s/3$ is $2400$~s at $t_s=7200$~s,
i.e.\ $10^\circ$ of longitude, which is exactly the period of the $m=36$
azimuthal harmonic on the $72$-column grid. That harmonic is sampled at its own
period, so instead of being averaged to zero it is passed at unit gain; the same
holds approximately for $m=35$--$40$, where the equal-spaced rule returns
$+0.99$ to $+0.84$ against true responses of $\pm0.07$. The worst error over
$m\le40$ at $t_s=7200$~s is $1.02$ --- larger than the response itself. This is a
quadrature alias, not a smoothing error: the rule is blind to the very structure
it is meant to average out, and data and operator generated with it are
self-consistent, so the mismatch is invisible in residuals and shows up only in
the recovered map.}

\rev{Table~\ref{tab:quad} reports $\max_{m\le40}|R_m^{(n_{\rm sub})}-R_m|$ for
Gauss--Legendre rules, computed by
{\tt sglsim.exposure\_quadrature\_convergence}. Convergence is monotonic in
$n_{\rm sub}$ at fixed $t_s$ and the GL rules never recover unit gain at an
aliased harmonic: GL-3 peaks at $0.88$ rather than $1.02$, and GL-16 is
\emph{spectrally converged} --- its worst error over all five dwell lengths is
$2.6\times10^{-13}$, i.e.\ machine precision even for the 7200~s dwells of
Arm~A, where the intra-dwell smear alone is $30^\circ$. GL-16 is therefore used
everywhere: in the data, in the forward operator, in the $\sigma_{\rm cl}$
calibration of Section~\ref{sec:ledger}, and in every reconstruction. The
practical consequence is that the $t_s=7200$~s ($M_p=4$) cadence points of the
original analysis were computed with an operator whose response error at the
highest retained harmonic is of order unity; they are superseded by the
Arm~A/Arm~B numbers of Section~\ref{sec:cadence}.}

\begin{table*}[t]
\centering
\caption{\rev{Maximum rotational-harmonic response error over $1\le m\le40$
against the exact sinc, by dwell length $t_s$. The top row is the equal-spaced
rule used in v2; the rest are Gauss--Legendre. From
{\tt results/quadrature.npz} (\texttt{eq3\_max\_err}, \texttt{gl\_max\_err}).}}
\label{tab:quad}
\footnotesize
\setlength{\tabcolsep}{3pt}
\begin{tabular}{lccccc}
\toprule
\rev{Rule} & \multicolumn{5}{c}{\rev{$\max_{m\le40}|R_m^{(n)}-R_m|$ by dwell $t_s$}}\\
\cmidrule(lr){2-6}
\rev{} & \rev{7200} & \rev{3600} & \rev{1800} & \rev{900} & \rev{450~s}\\
\midrule
\rev{equal-3 (v2)} & \rev{$1.0$}    & \rev{$1.3\times10^{-1}$} & \rev{$3.6\times10^{-2}$} & \rev{$2.4\times10^{-2}$} & \rev{$7.4\times10^{-3}$}\\
\rev{GL-3}  & \rev{$8.8\times10^{-1}$} & \rev{$2.7\times10^{-1}$} & \rev{$8.3\times10^{-3}$} & \rev{$1.5\times10^{-4}$} & \rev{$2.5\times10^{-6}$}\\
\rev{GL-5}  & \rev{$5.1\times10^{-1}$} & \rev{$3.5\times10^{-3}$} & \rev{$5.3\times10^{-6}$} & \rev{$5.8\times10^{-9}$} & \rev{$5.8\times10^{-12}$}\\
\rev{GL-8}  & \rev{$4.8\times10^{-3}$} & \rev{$2.4\times10^{-7}$} & \rev{$4.9\times10^{-12}$} & \rev{$4.4\times10^{-16}$} & \rev{$2.2\times10^{-16}$}\\
\rev{GL-16} & \rev{$2.6\times10^{-13}$} & \rev{$6.4\times10^{-16}$} & \rev{$6.4\times10^{-16}$} & \rev{$2.2\times10^{-16}$} & \rev{$2.2\times10^{-16}$}\\
\rev{GL-32} & \rev{$1.2\times10^{-15}$} & \rev{$1.1\times10^{-15}$} & \rev{$1.0\times10^{-15}$} & \rev{$4.4\times10^{-16}$} & \rev{$2.2\times10^{-16}$}\\
\bottomrule
\end{tabular}
\end{table*}

\rev{\subsection{What profiling actually removes: an information audit}
\label{sec:infaudit}
The rank/percentage framing of the v2 deflation paragraph cannot be repaired,
so we replaced it with a direct eigenvalue audit of the normal matrix at the
fiducial campaign (Table~\ref{tab:infaudit}, produced by
{\tt sglsim.profiling\_information\_audit}). Three findings.

First, profiling removes \textbf{no rank at all}. The un-profiled normal matrix
already has nullity 288, the free profile has nullity 288, and the anchored
profile has nullity 288. The null space is a property of the reconstruction
grid and the raster, not of the nuisance treatment: 144 basis columns are
\emph{identically zero} because they live in the two polar latitude rows of the
equirectangular grid, which a circular image-plane raster never samples (the
projected disk does not reach $|{\rm lat}|=90^\circ$ cell centres). The
remaining 144 null directions are the associated near-degenerate partners; the
smallest retained eigenvalue with $\lambda_{\rm eff}$ regularization is
$6.3\times10^{-6}$ against a maximum of $2.9\times10^{3}$.

Second, what profiling removes is the \emph{response along the mean-albedo
direction}, measured as the constant-mode gain $\|\mathsf{A}\mathbf{1}\|/\|\mathbf{1}\|$:
$913.7$ un-profiled, $18.79$ for the free profile ($2.06\%$ retained), and
$774.1$ for the anchored profile ($84.72\%$ retained). The free profile is
therefore not a mild correction --- it attenuates the disk-integrated albedo
by a factor of 49, which is why anchoring is mandatory rather than cosmetic.

Third, the identity $\mathsf{W}_\perp\mathsf{B}=0$ quoted in v2 is a trivial
property of any profile likelihood (the residual is orthogonal to the nuisance
space by construction); it is not evidence of a surface-mode degeneracy, and
an earlier version of this work read it as one.}

\begin{table}[t]
\centering
\caption{\rev{Information audit of slot-offset profiling at the fiducial
campaign ($n=64$, $N_{\rm sc}=16$, $M_p=16$, $n_s=2592$). Nullity is unchanged
by profiling; the constant-mode gain is what changes. \texttt{free} = offsets
projected out, \texttt{anchored} = Eq.~\ref{eq:anchor}.}}
\label{tab:infaudit}
\footnotesize
\setlength{\tabcolsep}{1.6pt}
\renewcommand{\arraystretch}{1.05}
\begin{tabular}{lccc}
\toprule
\rev{Quantity} & \rev{un-profiled} & \rev{free} & \rev{anchored} \\
\midrule
\rev{nullity} & \rev{288} & \rev{288} & \rev{288} \\
\rev{exactly-zero basis columns} & \rev{144} & \rev{144} & \rev{144} \\
\rev{constant-mode gain} & \rev{913.7} & \rev{18.79} & \rev{774.1} \\
\rev{fraction of gain retained} & \rev{1} & \rev{0.0206} & \rev{0.8472} \\
\rev{largest eigenvalue} & \rev{2940} & \rev{429.9} & \rev{1030} \\
\rev{$\lambda_{\rm eff}$} & \rev{0.00842} & \rev{0.00558} & \rev{0.00564} \\
\bottomrule
\end{tabular}
\end{table}

\subsection{Numerical precision verification}
\label{sec:precisionnum}
The normal matrix $\mathsf{A}=\mathsf{F}^{\sf T}\mathsf{C}^{-1}\mathsf{F}$
is accumulated in single precision (`float32`) and solved in double precision
(`float64`).
\rev{The v2 precision test was vacuous: it compared two double-precision
accumulations and reported their difference, which measures nothing. It is
replaced by a test that actually probes the single-precision path --- the same
design matrix, rounded to float32 before accumulation, versus the float64
reference --- at a reduced but honest problem size ($n=32$, 16{,}384 samples;
the fiducial $n=64$ design would need 2.8~GB in two float64 copies inside the
test harness), with the results in Table~\ref{tab:precision}.}
\rev{Three numbers matter. The float32 design matrix differs from the float64
reference by $3.1\times10^{-7}$ in Frobenius norm and the right-hand side by
$4.1\times10^{-7}$, i.e.\ at the float32 machine-precision level, and the same
quantities differ by only $1.0\times10^{-8}$ and $1.5\times10^{-8}$ from
rounding alone --- confirming the discrepancy is representational, not
algorithmic. Second, changing the \emph{summation order} of the blockwise
float64 accumulation perturbs $\mathsf{A}$ by $5.3\times10^{-16}$, which is the
number the v2 test was reporting. Third, and most relevant for the science:
solving the regularized normal equations versus a QR solve of the same whitened
design changes the recovered map by $1.7\times10^{-13}$ relative to the map
norm, and chunking the profiler in different orders changes the fiducial map by
$3.4\times10^{-13}$.}
\rev{The unregularized normal matrix is exactly singular (rank 2016 of 2592 at
the reduced size, nullity 288 at the fiducial size, Table~\ref{tab:infaudit}),
so $\kappa(\mathsf{A})$ is reported only after regularization, where
$\kappa=5.3\times10^{5}$ and the normal-equation residual is
$1.3\times10^{-15}$. Single-precision accumulation is therefore adequate for
this estimator, but the claim is now stated with its scope: it is a statement
about the linear algebra, verified at reduced size, not a general bound.}

\begin{table}[t]
\centering
\caption{\rev{Numerical precision audit (\texttt{sglsim.check\_numerical\_precision},
reduced size $n=32$, $N_{\rm sc}=16$, $M_p=16$).}}
\label{tab:precision}
\footnotesize
\setlength{\tabcolsep}{3pt}
\begin{tabular}{lc}
\toprule
\rev{Quantity measured} & \rev{Relative error} \\
\midrule
\rev{float32 vs float64 design $\mathsf{A}$ (Frobenius)} & \rev{$3.1\times10^{-7}$} \\
\rev{float32 vs float64 right-hand side $\mathbf{b}$} & \rev{$4.1\times10^{-7}$} \\
\rev{rounding-only contribution to $\mathsf{A}$} & \rev{$1.0\times10^{-8}$} \\
\rev{float64 summation-order change to $\mathsf{A}$} & \rev{$5.3\times10^{-16}$} \\
\rev{normal-equation vs QR solve (map)} & \rev{$1.7\times10^{-13}$} \\
\rev{profiler chunk-order change (map)} & \rev{$3.4\times10^{-13}$} \\
\rev{regularized condition number $\kappa(\mathsf{A})$} & \rev{$5.3\times10^{5}$} \\
\rev{unregularized rank / grid size} & \rev{$2016/2592$} \\
\bottomrule
\end{tabular}
\end{table}

\subsection{Baselines and statistical metrics}
\label{sec:baselines}
\rev{We compare TDI against three baselines, all evaluated on the same surface grid and with the same debiasing:
\begin{itemize}
\item {\bf B1 (phase-blind coadd):} the static limit of the same pipeline (a single rotational bin: the all-visit mean raster, Wiener-deconvolved and back-projected through the mean illumination; cf.\ \citealt{toth21}).
\item {\bf B2 (phase-registered coadd):} the phase-registered coadd baseline emulating \citet{tur26bench}: samples are binned in rotational phase, each bin is Wiener-deconvolved, and the bins are back-projected onto the sphere by bilinear interpolation weighted by illumination.
\item {\bf B3 (white-noise GLS):} a baseline generalized least-squares fit using the full time-domain operator $\mathsf{F}$ under an uncorrelated white-noise assumption ($\mathsf{C} = \mathrm{diag}(\sigma_k^2 + \sigma_{\rm cl}^2)$) without simultaneous-slot nuisance profiling (corresponding identically to row A of Table~\ref{tab:ablation}). This isolates the contribution of the time-dependent forward operator alone, in the absence of both nuisance profiling and temporal covariance modeling.
\end{itemize}
}

\rev{The bin count of B2 is not a free parameter and we fix it explicitly.
A pass visits a given image-plane pixel once, and consecutive dwell slots are
$t_s=1800$~s apart, i.e.\ $1/48$ of the rotation period: there are 48 distinct
sample phases per rotation. Coadding into only 8 phase bins averages over
$45^\circ$ of longitude, and on ideal cloud-free noiseless data that baseline
collapses to $r=0.07$ for reasons that have nothing to do with time-domain
inversion; $r=0.94$ is obtained at 16 bins, which is the best-tuned compromise
between rotational smear and the number of samples per bin ($M_p=16$ passes here, providing 48 sample phases per rotation). All B1/B2 numbers quoted below use $n_{\rm bins}=16$ for B2 and the
same Wiener regularization $K_w=3\times10^{-3}$; the naive 8-bin variant is
reported in Table~\ref{tab:clouds} only to quantify the artifact, so that the
competitor is not a straw man.}

\rev{\emph{Statistical reporting, geography conditioning, and error conventions.}
All synthetic campaigns in Sections~\ref{sec:rotation}--\ref{sec:lamsens} use a single fixed planetary geography ($\mathrm{TRUTH\_SEED}=7$, generating the synthetic continental albedo map shown in Fig.~\ref{fig:model}b). Between realization seeds (seeds 11--20), this planetary geography is held strictly constant while the photon shot noise (keyed to $\mathrm{seed}$) and the dynamic stochastic cloud field (keyed to $100+\mathrm{seed}$) are varied independently. Consequently, all uncertainties and confidence intervals reported in Tables~\ref{tab:clouds}, \ref{tab:cadence_ablation}, \ref{tab:ablation}, \ref{tab:geom}, \ref{tab:morph}, \ref{tab:wx}, and \ref{tab:sys} represent weather-and-photon sampling variability conditional on this single synthetic planet, not dispersion across diverse exoplanetary geomorphologies. The only true geomorphic scene diversity explored in this work is the three-morphology sweep of Section~\ref{sec:morph} (Table~\ref{tab:morph}).

Throughout this paper, the notation $x \pm y$ denotes the ensemble mean and standard error of the mean (SE, $\sigma/\sqrt{N_{\rm seeds}}$ with $N_{\rm seeds}=10$ unless stated otherwise), while standard deviation is explicitly labelled as SD and bootstrap intervals as 95\% CI (computed using 20{,}000 percentile resamples of paired differences; Appendix~\ref{sec:appendix_stats} (Table~\ref{tab:app_seeds}) tabulates all individual paired per-seed values and 95\% bootstrap intervals).
Because 10 seeds is a small ensemble, we report the exact attainable
significance floor alongside every $p$-value: for $n=10$ paired samples the
smallest attainable two-sided exact sign-test (binomial) $p$ is $2/2^{10}=1.95\times10^{-3}$, so no sign test on this ensemble can support $p<10^{-3}$. The paired $t$-test is a different
statistic and does reach far smaller $p$-values when the effect is consistent
in magnitude as well as sign; wherever a $p$ below the sign-test floor appears
in this paper it is a $t$-test $p$, and we label it as $p_t$. We also report a physical cloud-only null in which the surface is replaced by the mean and only
the weather is retained, whose 40-member ensemble limits attainable $p$ to
$1/(N_{\rm null}+1)=0.024$.}
\rev{Finally, SSIM and the land--ocean contrast ratio are computed on the
coarse $36\times72$ reconstruction grid with the same latitude mask as $r$, and
the harmonic decomposition of Section~\ref{sec:clouds} uses real spherical
harmonics with $\cos\vartheta$ quadrature weights, not a 2D FFT on an
equirectangular array (which is not a spherical-harmonic decomposition: the
grid cells are not equal-area and the azimuthal sampling is aliased at high
$|{\rm lat}|$).}

\section{Results}
\label{sec:results}

\subsection{Planetary rotation under characterized geometry}
\label{sec:rotation}
\rev{\emph{Headline result, cloud-free.} For a rotating, orbit-illuminated
cloud-free planet observed with the fiducial 87.33-day campaign and exactly known
spin geometry, the profiled time-domain estimator recovers Pearson
$r=0.9786\pm0.0002$ (standard error over 10 paired weather seeds),
$\mathrm{SSIM}=0.8482\pm0.0010$, normalized RMSE $0.0628\pm0.0003$, albedo bias
$-0.0006\pm0.0006$, and land--ocean contrast ratio $1.056\pm0.001$
(Fig.~\ref{fig:gallery}, top center; Table~\ref{tab:clouds}, first column). All
numbers in this section and in Table~\ref{tab:clouds} come from the $16$-node
exposure quadrature of Section~\ref{sec:quad}, the exact anchored nuisance
profiling of Section~\ref{sec:profiling}, and phase bins registered into
$16$ rotational bins; they therefore differ from the v2 values
($r=0.990$, $\mathrm{SSIM}=0.918$), which used the approximate deflation and a
coarser quadrature.}

\rev{To attribute what remains, we compare three reference configurations on the
\emph{same} data. (i) The phase-blind coadd (B1), which treats the campaign as a
single static exposure, collapses to $r=0.1155\pm0.0004$: with rotation
unmodeled the surface is essentially unrecoverable, and the $+0.863$ gap to TDI
is the contribution of the time-registered forward operator alone. (ii)
Registering the coadds into $16$ rotational phase bins (B2) recovers most of the
signal, $r=0.9361\pm0.0001$, but with a $-0.048$ albedo bias and a contrast
ratio of only $0.844$: binning averages within each $22.5^{\circ}$ phase window
and flattens the map. (iii) A white-noise GLS fit with the same time-dependent
operator reaches $r=0.9802\pm0.0002$, i.e.\ $0.0015$ \emph{above} the profiled
estimator (paired difference $-0.0015$, 0/10 seeds in favour of whitening,
$t=-22.6$).}

\rev{Finally, we evaluate a static ideal-data control with the planet \emph{frozen}: the same $64\times64$ raster, the
same $16$ revisits and the same $1800$~s dwell --- so every raster position
receives exactly the same $8$~h of integration and there is no weather --- but
all passes see one identical pose, so there is no rotational smearing to model
and no longitudinal phase coverage to solve. Scored on the same seeds using the B3 white-noise GLS estimator
(archive \texttt{static\_reference.npz}), this frozen configuration
yields $r=0.4074\pm0.0254$ (standard deviation across seeds
$0.0804$, per-seed range $0.284$--$0.523$), $\mathrm{SSIM}=0.3604$, and normalized
RMSE $0.4299$.

Measured against the rotating cloud-free headline ($r=0.9786\pm0.0002$), rotation provides a paired advantage of
$\Delta r = +0.5712\pm0.0253$ ($t=22.5$, $p_t=3.3\times10^{-9}$, 10/10 seeds; against the B3 white-noise rotating arm of $0.9802$, the paired gain is $+0.5728\pm0.0254$). The rotating cloudy campaign sits between them: a paired loss of $\Delta r = -0.6364\pm0.0083$ ($t=-76.6$, $p_t=5.2\times10^{-14}$, 10/10 seeds) relative to the rotating cloud-free baseline (and $-0.6379\pm0.0083$ relative to the B3 rotating arm). The physical mechanism is operator conditioning rather than information loss: with the planet frozen,
the $16$ passes of each raster position produce identical rows in
$\mathsf{F}$, so the static design lacks the longitudinal rank diversity that rotation naturally provides, leaving the reconstruction heavily reliant on the regularizer. A frozen planet is thus not an easier inverse problem, but a degenerate one; planetary rotation is structurally advantageous, and the $\approx0.636$ drop in fidelity is driven by stochastic weather.}

\rev{With no weather present, temporal whitening provides no benefit because there is no temporal error covariance to whiten; the high fidelity ($r=0.9786$) is delivered entirely by the time-registered forward operator. Scale-dependent correlation $r(\ell)$ (Section~\ref{sec:res}) shows band means above $0.99$ for all degrees up to the $\ell\approx18$ grid ceiling, confirming that no rotational band is lost inside the supported harmonic range; the residual $1-r\simeq0.021$ arises from finite aperture and regularization. Rotation is therefore tractable under characterized spin geometry, and the imaging error budget is overwhelmingly dominated by weather.}

\begin{figure*}[t]
\centering
\includegraphics[width=0.98\textwidth]{figures/fig_gallery.pdf}
\caption{\rev{Reconstruction gallery for the fiducial 87.33-day campaign
($n=64$, $M_p=16$, $t_s=1800$~s, 16-node exposure quadrature).
Top: truth surface albedo; TDI reconstruction of the rotating cloud-free planet
($r=0.979$, $\mathrm{SSIM}=0.848$); phase-blind coadd (B1) at realized cover
$0.56$, which is essentially featureless ($r=0.042$). Bottom, all at realized
cover $\simeq0.56$: $16$-bin phase-binned coadd (B2), white-noise GLS (B3), and
TDI. The gross continental dichotomy is recovered, but fine features are erased
by weather.}}
\label{fig:gallery}
\end{figure*}

\subsection{Cloud degradation and effective spatial resolution}
\label{sec:clouds}
\rev{Table~\ref{tab:clouds} and Figure~\ref{fig:ssimfc} document the collapse of
fidelity with cloud cover at fixed single-campaign resources. The cover
realized by the OU field differs from the nominal parameter, so both are
quoted: mean realized cover is $0.237$, $0.398$, $0.558$ and $0.710$ for the
nominal $f_c=0.25,0.40,0.55,0.70$ runs.}

\rev{Mean correlation falls from $0.9786$ (cloud-free) to $0.5966\pm0.0087$ at
$23.7\%$ realized cover, $0.4538\pm0.0074$ at $39.8\%$, $0.3422\pm0.0082$ at
$55.8\%$ and $0.2491\pm0.0105$ at $71.0\%$. SSIM falls faster: $0.288\pm0.005$
already at $23.7\%$ cover and $0.061\pm0.003$ at $71.0\%$, i.e.\ small-scale
structure is erased first, which is the expected signature of a noise process
whose power is broad in $\ell$ (Section~\ref{sec:res}).}

\rev{Two comparisons in Table~\ref{tab:clouds} qualify the headline claim.
(1)~At every cloudy cover the profiled estimator (TDI) beats the white-noise GLS
fit (B3) by $\Delta r=+0.069$ to $+0.039$ (10/10 paired seeds in favour,
$t=9.3$ to $7.4$). However, as demonstrated by the component ablation in Section~\ref{sec:ablation} (Table~\ref{tab:ablation}), this performance gap is driven by the elimination of the simultaneous-slot cloud nuisance rather than by temporal covariance modeling: exact nuisance profiling alone recovers $\Delta r = +0.0568\pm0.0046$ over white GLS without profiling (Table~\ref{tab:ablation}, row D vs.\ row A), whereas adding OU temporal covariance to exact profiling slightly reduces correlation at this cadence ($\Delta r = -0.0032\pm0.0019$, row F vs.\ row D). B3 corresponds identically to Table~\ref{tab:ablation} row A and therefore differs from TDI in two factors simultaneously (slot handling and covariance model); Table~\ref{tab:ablation} provides the matched single-factor contrast.
(2)~At $23.7\%$ cover the $16$-bin phase-binned
coadd (B2) actually scores \emph{higher} than TDI, $r=0.6097$ vs.\ $0.5966$
(paired $\Delta r=-0.013$, 1/10 seeds in favour of TDI, $t=-2.62$), and the
gap only reverses at higher cover ($+0.006$ at $39.8\%$, $+0.025$ at
$55.8\%$, $+0.046$ at $71.0\%$, $t=0.79$ to $6.14$). B2 is a
non-parametric estimator: it makes no assumption about cloud statistics, so
at low cover --- where the cloud nuisance is small --- its binning variance
costs less than TDI's model error. The claim we can defend is therefore not
``TDI is always better than phase-binning,'' but ``TDI's advantage grows
with cover, and at high cover B2 additionally carries a bias that TDI does
not'' ($+0.065$ vs.\ $+0.019$ albedo bias and contrast ratio $0.786$ vs.\
$1.024$ at $71\%$ cover).}

\begin{table}[t]
\centering
\caption{\rev{Map fidelity vs.\ cloud cover for the fiducial campaign
($n=64$, $M_p=16$, $t_s=1800$~s, 16-node quadrature, exact anchored
profiling). Rows are means $\pm$ standard error over 10 paired weather
seeds; the B2/B1 rows use 16, 8 and 1 rotational phase bins respectively.
White-noise GLS (B3) corresponds identically to row A of Table~\ref{tab:ablation}
(no slot profiling, white noise); Table~\ref{tab:ablation} isolates the single-factor matched contrast.}}
\label{tab:clouds}
\footnotesize
\setlength{\tabcolsep}{2.2pt}
\resizebox{\columnwidth}{!}{%%
\begin{tabular}{lccccc}
\toprule
\rev{Method} & \rev{$f_c{=}0$} & \rev{0.24} & \rev{0.40} & \rev{0.56} & \rev{0.71}\\
\midrule
\multicolumn{6}{c}{\rev{Pearson $r$}}\\
\rev{TDI (profiled)}   & \rev{0.9786} & \rev{0.5966} & \rev{0.4538} & \rev{0.3422} & \rev{0.2491}\\
\rev{}                 & \rev{$\pm.0002$} & \rev{$\pm.0087$} & \rev{$\pm.0074$} & \rev{$\pm.0082$} & \rev{$\pm.0105$}\\
\rev{White GLS (B3)}   & \rev{0.9802} & \rev{0.5274} & \rev{0.3863} & \rev{0.2886} & \rev{0.2107}\\
\rev{}                 & \rev{$\pm.0002$} & \rev{$\pm.0085$} & \rev{$\pm.0080$} & \rev{$\pm.0097$} & \rev{$\pm.0107$}\\
\rev{Phase-binned (B2)}& \rev{0.9361} & \rev{0.6097} & \rev{0.4479} & \rev{0.3175} & \rev{0.2028}\\
\rev{}                 & \rev{$\pm.0000$} & \rev{$\pm.0083$} & \rev{$\pm.0121$} & \rev{$\pm.0158$} & \rev{$\pm.0157$}\\
\rev{Phase-blind (B1)} & \rev{0.1155} & \rev{0.0782} & \rev{0.0567} & \rev{0.0417} & \rev{0.0328}\\
\rev{}                 & \rev{$\pm.0004$} & \rev{$\pm.0145$} & \rev{$\pm.0126$} & \rev{$\pm.0086$} & \rev{$\pm.0058$}\\
\rev{B2, 8 bins}       & \rev{0.0700} & \rev{0.0136} & \rev{$-$0.0140} & \rev{$-$0.0309} & \rev{$-$0.0416}\\
\rev{}                 & \rev{$\pm.0005$} & \rev{$\pm.0193$} & \rev{$\pm.0222$} & \rev{$\pm.0246$} & \rev{$\pm.0224$}\\
\midrule
\multicolumn{6}{c}{\rev{SSIM}}\\
\rev{TDI (profiled)}   & \rev{0.8482} & \rev{0.2876} & \rev{0.1748} & \rev{0.1060} & \rev{0.0606}\\
\rev{White GLS (B3)}   & \rev{0.8568} & \rev{0.2492} & \rev{0.1447} & \rev{0.0853} & \rev{0.0472}\\
\rev{Phase-binned (B2)}& \rev{0.7947} & \rev{0.3515} & \rev{0.2223} & \rev{0.1307} & \rev{0.0627}\\
\rev{Phase-blind (B1)} & \rev{0.0831} & \rev{0.0376} & \rev{0.0143} & \rev{0.0028} & \rev{$-$0.0011}\\
\midrule
\multicolumn{6}{c}{\rev{albedo bias / contrast ratio}}\\
\rev{TDI (profiled)}   & \rev{$-$0.001/1.056} & \rev{$-$0.017/1.103} & \rev{$-$0.009/1.107} & \rev{$+$0.004/1.071} & \rev{$+$0.019/1.024}\\
\rev{Phase-binned (B2)}& \rev{$-$0.048/0.844} & \rev{$-$0.057/0.862} & \rev{$-$0.034/0.851} & \rev{$+$0.003/0.826} & \rev{$+$0.065/0.786}\\
\bottomrule
\end{tabular}
}
\end{table}

\begin{figure*}[t]
\centering
\includegraphics[width=0.95\textwidth]{figures/fig_ssim_fc.pdf}
\caption{\rev{Reconstruction fidelity vs.\ realized cloud cover across 10
paired seeds; bars are $\pm1$ standard error. The $16$-bin phase-binned
coadd is competitive below $\sim40\%$ cover and falls behind as weather
dominates (Table~\ref{tab:clouds}). Clouds, not rotation, control image
quality at fixed observing resources.}}
\label{fig:ssimfc}
\end{figure*}

\subsection{\rev{Effective spatial resolution and the land--ocean test}}
\label{sec:res}
\rev{Scale-dependent correlation.} Decomposing the maps in real spherical harmonics
with $\cos\vartheta$ quadrature weights over degrees $\ell\in[1,20]$
(Fig.~\ref{fig:robust}c) gives, over 10 seeds:
\rev{$r(\ell)=1.000\pm0.001$ averaged over $\ell=1$--$4$,
$0.9995$ over $5$--$8$, $0.9988$ over $9$--$12$, $0.9942$ over $13$--$16$ and
$0.9924$ over $17$--$20$ for the cloud-free campaign; the cloudy campaign
($55.8\%$ realized cover) gives $0.700$, $0.594$, $0.418$, $0.244$ and $0.263$
in the same bands.}
\rev{The v2 statement that the cloud-free curve ``extends past $\ell=20$''
cannot be tested on this grid: $36\times72$ pixels cannot represent degrees
above $\ell\approx18$, and the discrete-basis orthonormality error of the
quadrature itself is $0.018$ at $\ell=20$, so we quote $\ell\le20$ as the
measurement ceiling and report the ceiling error alongside. The cloud-free
result should therefore be read as ``no rotational band within the supported
range is lost,'' not as a resolution claim.}
\rev{For the cloudy case the curve is not monotone at the seed level (the
per-degree standard error over the ten seeds is $0.02$--$0.12$), so a single
crossing degree
$\ell_{\rm eff}$ is fragile: the first degree at which the raw per-seed curve
dips below $0.5$ is $\ell=1$ for five seeds and $\ell\in\{3,7,9\}$ for the
other five. The defensible
statement is a banded one: a five-point running mean of $r(\ell)$ crosses
$0.5$ at $\ell\simeq9$ (half-wavelength $\pi R_\oplus/\ell\approx2200$~km for
an Earth-radius planet), only $5/10$ seeds have $r(\ell)\ge0.5$ at
$\ell=7$, and beyond that the seed-mean $r(\ell)$ never recovers: for every
degree $\ell\ge10$ the mean sits below $0.55$ and for every $\ell\ge15$ below
$0.41$ (per-degree SE $0.03$--$0.06$ for $\ell\ge15$). Individual seeds still
clear $0.5$ at $\ell=16$ ($1/10$) and $\ell=20$ ($3/10$, mean $0.40$), which is the
\emph{per-degree} statement we can defend; v2's stronger ``no degree
$\ell\ge15$ reaches $0.5$ in more than $10\%$ of seeds'' is not reproducible on
the archived curves and is withdrawn. Clouds thus leave the hemispheric dichotomy and the largest
continents recoverable, and nothing at scales below $\sim2000$~km.}

\rev{Land--ocean detection.} We compute
$d'=(\mu_{\rm land}-\mu_{\rm ocean})/\sigma_{\rm pooled}$ within
$|\vartheta|\le60^\circ$, using the true land labels. \rev{Unlike the $r$
tables, the $d'$ values below are quoted as mean $\pm$ the per-seed
\emph{standard deviation} over the 10 seeds, because the quantity of interest is
the spread of the detection statistic across weather realizations; divide by
$\sqrt{10}$ for the corresponding standard error ($0.081\to0.026$).} The truth
field itself
gives $d'=3.97$. Cloud-free TDI gives $d'=3.58\pm0.03$ (SE $0.009$), the white
GLS fit $3.61\pm0.02$ (SE $0.007$), and B2 $3.04\pm0.005$ (SE $0.001$); at
$55.8\%$ realized cover TDI gives
$d'=0.732\pm0.081$ and B2 $0.674\pm0.131$. Significance is bounded by the
physical cloud-only null --- 40 realizations in which the surface is replaced
by its mean and only weather is retained --- which has
$\bar{d}'=0.001$, $\mathrm{sd}=0.090$ and 95th percentile $0.114$; the cloudy
TDI value exceeds every null member, but a 40-member ensemble cannot
support a $p$ value below $1/(N_{\rm null}+1)=0.024$. \rev{We therefore report
$p\le0.024$ (rank-based, one-sided) for the cloudy detection rather than the
v2 ``$p<0.02$ against 100 surrogate nulls,'' whose surrogate was a
phase-scrambled equirectangular array --- a geometry transform, not a
photometric null. A label-permutation null on the same reconstructions is
reported with the robustness suite in Section~\ref{sec:robust}.}

\rev{Quality is also a function of latitude, and the two effects do not point
the same way. Scoring the archived fiducial reconstructions band by band over
$|\vartheta|$ in $30^{\circ}$ steps (archive {\tt latband.npz}, 10 seeds) pairs
each band with the illumination it actually receives --- the campaign-mean
Lambert factor over its disk pixels, and the campaign-mean column energy
$\Vert F_{:,i}\Vert^{2}$, which is the information the operator delivers there.}

\begin{table}[t]
\centering
\caption{\rev{Latitude- and illumination-resolved quality for the fiducial
campaigns (10 seeds, $\pm1$~SE). ``Dayside'' is the campaign-mean Lambert
factor over the band's disk pixels; ``col.\ energy'' is the campaign-mean column
energy $\Vert F_{:,i}\Vert^{2}$ of the operator; ``dead'' is the fraction of the
band's $432$ grid cells whose column energy is below $10^{-3}$ of the maximum;
``gain'' is $\mathrm{sd}(\text{recon})/\mathrm{sd}(\text{truth})$ in the band,
which Pearson $r$ is blind to.}}
\label{tab:lat}
\footnotesize
\setlength{\tabcolsep}{2.6pt}
\renewcommand{\arraystretch}{1.05}
\resizebox{\columnwidth}{!}{%
\begin{tabular}{lccccccc}
\toprule
\rev{Latitude} & \rev{disk pix} & \rev{dayside} & \rev{col.\ energy} &
\rev{dead} & \rev{$r$ (cloudy)} & \rev{$r$ (cloud-free)} & \rev{gain (cloudy)} \\
\midrule
\rev{$-90$--$-60^{\circ}$} & \rev{86}  & \rev{0.258} & \rev{0.33}  & \rev{0.500} & \rev{$0.451\pm0.038$} & \rev{$0.9335\pm0.0027$} & \rev{1.95} \\
\rev{$-60$--$-30^{\circ}$} & \rev{546} & \rev{0.524} & \rev{6.61}  & \rev{0}     & \rev{$0.300\pm0.020$} & \rev{$0.9655\pm0.0008$} & \rev{3.89} \\
\rev{$-30$--$0^{\circ}$}   & \rev{982} & \rev{0.691} & \rev{21.43} & \rev{0}     & \rev{$0.353\pm0.015$} & \rev{$0.9814\pm0.0002$} & \rev{3.49} \\
\rev{$0$--$30^{\circ}$}    & \rev{982} & \rev{0.691} & \rev{21.43} & \rev{0}     & \rev{$0.390\pm0.025$} & \rev{$0.9863\pm0.0001$} & \rev{2.77} \\
\rev{$30$--$60^{\circ}$}   & \rev{546} & \rev{0.524} & \rev{6.60}  & \rev{0}     & \rev{$0.335\pm0.022$} & \rev{$0.9742\pm0.0006$} & \rev{2.89} \\
\rev{$60$--$90^{\circ}$}   & \rev{86}  & \rev{0.258} & \rev{0.33}  & \rev{0.486} & \rev{$0.547\pm0.038$} & \rev{$0.9358\pm0.0025$} & \rev{1.68} \\
\bottomrule
\end{tabular}%
}
\end{table}

\rev{Table~\ref{tab:lat} shows the ordering the referee expected and the one
they did not. Expected: band quality is monotone in the information the operator
actually delivers there. Across the well-sampled globe the two bands carrying
$21.4$ of column energy and $982$ disk pixels each score $0.353$ and $0.390$
(cloudy) and $0.9814$ and $0.9863$ (cloud-free), while the two bands carrying
$6.6$ score $0.300$ and $0.335$, and $0.9655$ and $0.9742$ --- the same ordering
in both campaigns, with no band out of place. Unexpected: the two polar bands score
\emph{highest} in the cloudy case ($0.451$ and $0.547$). That is an artefact of
the metric, not a recovery claim, and the table gives the two diagnostics that
expose it. The polar bands contain a half-illuminated ice cap against a dark
ocean, so their truth standard deviation is the largest of any band
($0.148$--$0.160$ against $0.097$--$0.137$ mid-latitude) while their structure is
dominated by the low-degree cap edge. A reconstruction that
gets nothing but the smooth cap still correlates well with truth. The amplitude
column confirms it --- the polar gain is $1.68$--$1.95$ whereas the
mid-latitude bands are inflated $2.8$--$3.9\times$ --- and restricting the polar
correlation to the cells the operator actually samples (all grid cells whose
campaign-mean column energy exceeds $10^{-3}$ of the maximum, i.e.\ the
complement of the dead fraction in the last column) raises it further, to
$0.518\pm0.052$ and $0.652\pm0.034$. Half of each polar band is dead: the circular $n=64$
raster never samples $|\vartheta|\gtrsim85^{\circ}$ (the
$144$ exactly-zero basis columns of Section~\ref{sec:infaudit}), so those cells
are set by the regularizer alone and contribute truth-like structure to nothing
but the denominator of $r$. Cloud-free the same geometry gives the clean
illumination monotonic trend, $0.9863$ at the equator to $0.934$ at the poles,
with the polar deficit already concentrated in the dead cells
($0.978$--$0.974$ restricted to sampled cells). The band offsets are
symmetric-to-within-noise --- $|\overline{\Delta a}| \le 0.019$ cloudy and
$\le0.005$ cloud-free --- so the latitude dependence is a dependence of
\emph{correlation}, not a latitudinal albedo bias. We quote the band-resolved
$r$ as a coverage diagnostic and explicitly not as a claim of polar imaging.}
\subsection{Component-wise ablation of the inversion pipeline}
\label{sec:ablation}
\rev{Table~\ref{tab:ablation} ablates the estimator itself across 10 paired
seeds at nominal $f_c=0.55$ ($0.558$ realized cover), holding the data, the
operator, the grid and the regularizer fixed and changing only how the
dwell-slot nuisance and the temporal covariance are handled. Row C isolates the approximate deflation heuristic under the OU covariance model, whereas earlier exploratory runs altered forward-model choices and deflation heuristics together without isolating individual mechanisms.}

\begin{table}[t]
\centering
\caption{\rev{Estimator-component ablation at $f_c=0.55$ (10 paired seeds;
means $\pm$ standard error). ``Profiling'' is the exact anchored nuisance
elimination of Section~\ref{sec:profiling}; ``deflation'' is the
heuristic approximation to it; ``whitening'' is the OU temporal covariance. The $n/10$
entry is the number of seeds whose paired difference carries the sign of the
mean; $p_t$ is the paired $t$-test $p$ (the exact sign-test floor is $1.95\times10^{-3}$ at $n=10$).}}
\label{tab:ablation}
\footnotesize
\setlength{\tabcolsep}{2.0pt}
\resizebox{\columnwidth}{!}{%%
\begin{tabular}{lccc}
\toprule
\rev{Configuration} & \rev{Pearson $r$} & \rev{SSIM} & \rev{NRMSE}\\
\midrule
\rev{A: white, no slot handling}   & \rev{$0.2886\pm0.0097$} & \rev{$0.0853\pm0.0031$} & \rev{$0.966$}\\
\rev{B: OU whitened, no slot handling} & \rev{$0.2863\pm0.0090$} & \rev{$0.0825\pm0.0031$} & \rev{$0.974$}\\
\rev{C: approximate deflation, OU} & \rev{$0.3422\pm0.0082$} & \rev{$0.1058\pm0.0055$} & \rev{$0.814$}\\
\rev{D: exact profiling, white, anchored} & \rev{$0.3454\pm0.0082$} & \rev{$0.1095\pm0.0051$} & \rev{$0.809$}\\
\rev{E: exact profiling, white, free} & \rev{$0.3453\pm0.0082$} & \rev{$0.1093\pm0.0054$} & \rev{$0.809$}\\
\rev{F: full TDI (exact profiling + OU)} & \rev{$0.3422\pm0.0082$} & \rev{$0.1060\pm0.0053$} & \rev{$0.814$}\\
\midrule
\rev{F$-$A ($\Delta r$)} & \multicolumn{2}{c}{\rev{$+0.0536\pm0.0061$ ($t=8.8$, 10/10)}} & \rev{$p_t=1.1\times10^{-5}$}\\
\rev{D$-$A ($\Delta r$)} & \multicolumn{2}{c}{\rev{$+0.0568\pm0.0046$ ($t=12.4$, 10/10)}} & \rev{$p_t=5.9\times10^{-7}$}\\
\rev{F$-$C ($\Delta r$, matched)} & \multicolumn{2}{c}{\rev{$+4.1\times10^{-5}\pm3.0\times10^{-5}$ ($t=1.38$, 7/10 in sign)}} & \rev{$p_t=0.20$}\\
\rev{D$-$C ($\Delta r$)} & \multicolumn{2}{c}{\rev{$+0.0032\pm0.0019$ ($t=1.7$, 6/10 in sign)}} & \rev{$p_t=0.13$}\\
\rev{F$-$D ($\Delta r$)} & \multicolumn{2}{c}{\rev{$-0.0032\pm0.0019$ ($t=-1.7$, 6/10 in sign)}} & \rev{$p_t=0.13$}\\
\bottomrule
\end{tabular}
}
\end{table}

\rev{Three conclusions follow from this matched factorial decomposition.}

\rev{(i) The gain is nuisance handling, and it is real. Exact profiling lifts
$r$ by $+0.057\pm0.005$ over white GLS without slot handling ($D - A$, $t=12.4$, $p_t=5.9\times10^{-7}$, 10/10 seeds) and $+0.054\pm0.006$ over white GLS when combined with OU covariance ($F - A$, $t=8.8$, $p_t=1.1\times10^{-5}$), with
$\Delta\mathrm{SSIM}=+0.021\pm0.004$ ($t=5.6$). Eliminating the per-slot nuisance mean recovers essentially the entire algorithmic gap between TDI and standard inversion.}

\rev{(ii) Exact profiling vs.\ approximate deflation is a matched contrast ($F - C$, holding OU covariance fixed). Under this matched comparison, exact profiling yields $\Delta r = +4.1\times10^{-5}\pm3.0\times10^{-5}$ ($t=1.38$, $p_t=0.20$) and a relative $L_2$ map difference of only $0.83\%\pm0.22\%$ in norm. The earlier deflation heuristic was therefore quantitatively very close to the exact profile likelihood solution in reconstruction space. We adopt the exact treatment because it is mathematically derived from the constrained likelihood (Eqs.~\ref{eq:nuisance}--\ref{eq:anchordef}), not because it yields a statistically resolvable gain in $r$ over the heuristic.}

\rev{(iii) Temporal whitening contributes essentially nothing at the fiducial cadence, but accounts for the map difference previously misattributed to deflation. Comparing $F$ against $D$ (OU vs.\ white under exact profiling) gives $F-D = -0.0032\pm0.0019$ ($t=-1.7$, $p_t=0.13$), and comparing $B$ against $A$ gives $B-A = -0.0023\pm0.0015$. However, the relative $L_2$ norm difference between $D$ (white) and $F$ (OU) is $9.22\%$, and between $D$ and $C$ (OU) is $9.29\%$ --- confirming that the $\approx9.3\%$ map change is caused by the switch between temporal covariance weighting and white-noise weighting, not by the slot-deflation approximation. With 16 revisits spaced $5.47$~d and $\tau_{\rm cloud}=4$~d, consecutive passes are largely decorrelated ($\rho_1=0.255$; Table~\ref{tab:neff}), so OU whitening provides little additional correlation and incurs a minor variance penalty.}

\rev{Finally, the anchored (D) and free (E) profiles yield identical Pearson correlation $r = 0.3453$ to four decimal places of display precision, but form a statistically significant paired contrast $E - D = -9.6\times10^{-5}\pm2.8\times10^{-5}$ ($t = -3.39$, $p_t = 0.008$), with a map norm difference of $0.98\%\pm0.27\%$. Their zero-mode conditioning differs by a factor of 41 (Section~\ref{sec:infaudit}): the anchor identifies the absolute albedo scale without altering spatial correlation.}

\subsection{Isolating the causal mechanism of the cadence law}
\label{sec:cadence}
\rev{Figure~\ref{fig:cadence} and Table~\ref{tab:cadence_ablation} report the
cadence law under the two resource ledgers of Section~\ref{sec:ledger}, at
nominal $f_c=0.55$. Each cell is a mean $\pm$ standard error over the six
paired seeds of the campaign study; the four estimator columns are the
profiled TDI, the white-noise GLS fit, the 16-bin phase-binned coadd, and the
phase-blind coadd. The v2 table quoted $r_{\rm inst}$, $r_{\rm smear}$,
$r_{\rm cf}$ and $r_{\rm iid}$ columns that the shipped pipeline never
produced; they have been replaced by controls that are in the archive
(Section~\ref{sec:cadctl}).}

\begin{table}[t]
\centering
\caption{\rev{Cadence law per resource arm (nominal $f_c=0.55$, 6 paired
seeds). Arm A holds $M_pt_s=8$~h and lets the wall clock exceed 90~d;
arm B holds 89.87~d and lets $t_s$ fall. $\eta$ is the duty cycle under
45~s/dwell overhead.}}
\label{tab:cadence_ablation}
\footnotesize
\setlength{\tabcolsep}{2.6pt}
\resizebox{\columnwidth}{!}{%%
\begin{tabular}{ccccccccc}
\toprule
\rev{Arm} & \rev{$M_p$} & \rev{$t_s$} & \rev{$\eta$} & \rev{wall} & \rev{$r_{\rm TDI}$} & \rev{$r_{\rm white}$} & \rev{$r_{\rm B2}$} & \rev{$r_{\rm B1}$}\\
\midrule
\rev{A} & \rev{4} & \rev{7200} & \rev{0.994} & \rev{85.73} & \rev{$0.049\pm0.005$} & \rev{$0.024\pm0.003$} & \rev{$0.063\pm0.011$} & \rev{$0.003\pm0.020$}\\
\rev{A} & \rev{8} & \rev{3600} & \rev{0.988} & \rev{86.27} & \rev{$0.173\pm0.013$} & \rev{$0.144\pm0.011$} & \rev{$0.068\pm0.011$} & \rev{$0.063\pm0.019$}\\
\rev{A} & \rev{16} & \rev{1800} & \rev{0.976} & \rev{87.33} & \rev{$0.346\pm0.010$} & \rev{$0.291\pm0.009$} & \rev{$0.321\pm0.019$} & \rev{$0.031\pm0.010$}\\
\rev{A} & \rev{32} & \rev{900} & \rev{0.952} & \rev{89.47} & \rev{$0.497\pm0.014$} & \rev{$0.475\pm0.012$} & \rev{$0.074\pm0.004$} & \rev{$-0.019\pm0.003$}\\
\rev{A} & \rev{64} & \rev{450} & \rev{0.909} & \rev{\bfseries 93.73} & \rev{$0.571\pm0.010$} & \rev{$0.573\pm0.013$} & \rev{$0.503\pm0.013$} & \rev{$0.179\pm0.002$}\\
\midrule
\rev{B} & \rev{4} & \rev{7548.8} & \rev{0.994} & \rev{89.87} & \rev{$0.049\pm0.006$} & \rev{$0.031\pm0.005$} & \rev{$-0.130\pm0.005$} & \rev{$-0.064\pm0.017$}\\
\rev{B} & \rev{8} & \rev{3751.9} & \rev{0.988} & \rev{89.87} & \rev{$0.164\pm0.014$} & \rev{$0.140\pm0.013$} & \rev{$-0.134\pm0.021$} & \rev{$-0.024\pm0.019$}\\
\rev{B} & \rev{16} & \rev{1853.4} & \rev{0.976} & \rev{89.87} & \rev{$0.325\pm0.017$} & \rev{$0.267\pm0.015$} & \rev{$-0.231\pm0.022$} & \rev{$0.015\pm0.014$}\\
\rev{B} & \rev{32} & \rev{904.2} & \rev{0.953} & \rev{89.87} & \rev{$0.498\pm0.008$} & \rev{$0.489\pm0.010$} & \rev{$-0.140\pm0.058$} & \rev{$-0.025\pm0.006$}\\
\rev{B} & \rev{64} & \rev{429.6} & \rev{0.905} & \rev{89.87} & \rev{$0.576\pm0.012$} & \rev{$0.624\pm0.008$} & \rev{$0.454\pm0.018$} & \rev{$0.076\pm0.002$}\\
\bottomrule
\end{tabular}
}
\end{table}

\rev{\subsection{What the cadence law does and does not say}
\label{sec:cadctl}
\emph{The trend is monotone, and it is dominated by weather sampling rather
than by conditioning.} In arm A, $r$ rises $0.049\to0.571$ as $M_p$ goes $4\to64$, and every
step is unanimous across the six paired seeds ($t=8.7$ to $21.9$; marginal
gaps $+0.124,+0.173,+0.151,+0.075$). The gain saturates: the last doubling buys
less than half of the previous one. The cloud-free control of
Section~\ref{sec:rotation} shows that the collapse from $0.979$ to $0.346$ at the
fiducial cadence is entirely the weather term, and more passes help because each
pass is a fresh draw of it.}

\rev{\emph{The two controls, now actually measured.} The cloud-free and
per-pass-independent controls are computed on the same campaign objects, seeds
and estimator as the main sweep (six paired seeds, both arms). Cloud-free $r$:
$0.773\pm0.008$, $0.934\pm0.001$, and then $0.979$, $0.984$, $0.987$ to within
$\pm0.0002$ for $M_p=16,32,64$ in arm A, with arm B within $0.04$ of each.
Per-pass independent weather: $0.051$, $0.165$, $0.342$, $0.494$, $0.583$,
against the correlated-OU values $0.049$, $0.173$, $0.346$, $0.497$, $0.571$ in
the same order. Two conclusions follow, and they point in opposite directions.
(i) \emph{Conditioning contributes, but it saturates}: the cloud-free curve
still rises with $M_p$, by $0.214$ over $4\to64$, so part of the v2-era trend is
coverage and not weather averaging. It is, however, a minority of the effect and
it is front-loaded: it gains $0.16$ in the first doubling ($M_p=4\to8$, where a
raster position receives only four dwells) and only $0.008$ from $M_p=16\to64$,
while the cloudy curve is still climbing steeply there ($0.346\to0.571$,
$+0.225$). (ii) \emph{The law does not need OU structure}: the iid and OU
reconstructions agree at every cadence --- paired differences $+0.002$, $-0.008$,
$-0.004$, $-0.003$, $+0.012$ in arm A and $+0.001$, $-0.010$, $+0.014$, $-0.026$,
$+0.002$ in arm B, every $p_t\ge0.14$, none approaching the exact sign-flip
floor of $0.031$ attainable with six pairs. What revisits buy is independent
samples of the atmosphere, whatever its temporal correlation function turns out
to be; the OU model determines \emph{how many} of the $M_p$ looks are effectively
independent (Table~\ref{tab:neff}), it is not the reason that more looks help.
The absolute cost of weather at fixed cadence, $\Delta r_{\rm cf-OU}$, falls
monotonically from $+0.724\pm0.010$ at $M_p=4$ to $+0.416\pm0.011$ at $M_p=64$
(arm A; both arms significant at $p_t<4\times10^{-7}$), which is the same
statement in the other direction: the cadence law is about averaging the weather
down, not about the weather getting worse.}

\rev{\emph{The two arms agree within the paired standard errors.} Pairing arm A
against arm B at matched $M_p$ gives $\Delta r$ of $-0.001,+0.008,+0.021,
-0.002,-0.005$ with paired standard errors $0.005$--$0.018$ and $p\ge0.058$ in
every case. The wall-clock cost of doubling the revisit count is therefore paid
out of the photon budget with no measurable loss over the range studied: the
law is about \emph{independent looks}, not about integration time per look, so
long as the dwell stays long enough to beat the overhead (duty cycle
$0.994\to0.905$ across the arms).}

\rev{\emph{Whitening loses its edge at the fastest cadence.} The profiled
estimator beats the white fit at every $M_p\le32$ in both arms ($+0.025$ to $+0.058$ in Arm A; $+0.009$ to $+0.058$ in Arm B), but at
$M_p=64$ the ordering reverses in both arms: $-0.002\pm0.009$ in arm A and
$-0.047\pm0.009$ (0/6 seeds) in arm B. At 1.47~d spacing the OU correlation
$\rho_1=0.69$ (Table~\ref{tab:neff}) is real, but the estimator's assumed
$\sigma_{\rm cl}$ is calibrated at $t_s=1800$~s, and the whitening then
over-weights the residual structure. The safe reading is that temporal
whitening pays between roughly 8 and 32 revisits, and the fiducial 16 is not
the optimum of the whitened fit.}

\rev{\emph{The phase-binned baseline is partly a coverage artifact, and we can
show it.} B2 does not merely degrade in arm B, it goes strongly
\emph{negative} ($-0.23$ at $M_p=16$), which is not a statement about TDI. The
cause is geometric. Registering samples into 16 rotational phase bins does not
leave each bin with a full raster: counting, per bin, how many of the 4096
image-plane pixels receive no sample at all gives an average of 256 empty
pixels per bin at the fiducial configuration (arm A, $M_p=16$), but 2048 at
$M_p=16$ in arm B and 3584 at $M_p=4$ --- the phase pattern of a
boustrophedon raster does not tile the bins uniformly, and the two arms place
different dwells at different phases. \texttt{reconstruct\_phasebin} fills an
empty pixel with the bin's raster mean and then back-projects, so a
coverage-starved bin contributes a smoothed, self-referential map, exactly the
artifact flagged for the 8-bin variant in Section~\ref{sec:baselines} but an
order of magnitude worse. The B2 column is therefore reported for
completeness and is \emph{not} used to argue the cadence law; the comparison we
stand behind is TDI against the white-noise GLS fit, which shares the operator
and has no rasterization step.}

\rev{\emph{Overheads bound the law from below, not above.} With
$T_{\rm oh}=45$~s fixed, $\eta$ falls from $0.994$ to $0.905$; extrapolating
the same ledger to $t_s=100$~s costs 31\% of the photons and to $t_s=45$~s
costs half of them, so the ``ultrafast cadence'' limit is where overheads bite.
The v2 sentence claiming overhead losses ``offset statistical gains'' at
$t_s\lesssim100$~s was not measured --- no configuration below 429.6~s is in
the archive --- and is deleted.}

\rev{\begin{table}[t]
\centering
\caption{\rev{Effective number of independent looks,
$N_{\rm eff}=M_p^2/(M_p+2\sum_{k\ge1}(M_p-k)\rho_k)$ with
$\rho_k=e^{-k\Delta t/\tau_{\rm cloud}}$ and per-pass revisit spacing
$\Delta t=256(t_s+45)/86{,}400$~d (tabulated for Arm A, with $t_s=7200, 3600, 1800, 900, 450$~s; for Arm B, the corresponding dwells yield closely matching values, e.g., $N_{\rm eff}=11.64$ at $M_p=64$ for $\tau=4$~d).}}
\label{tab:neff}
\footnotesize
\setlength{\tabcolsep}{4pt}
\resizebox{\columnwidth}{!}{%%
\begin{tabular}{lccccc}
\toprule
\rev{$\tau_{\rm cloud}$} & \rev{$M_p=4$} & \rev{8} & \rev{16} & \rev{32} & \rev{64}\\
\midrule
\rev{spacing (d)}          & \rev{21.47} & \rev{10.80} & \rev{5.47} & \rev{2.80} & \rev{1.47}\\
\rev{$\rho_1$ (lag-1)}     & \rev{0.005} & \rev{0.067} & \rev{0.255} & \rev{0.497} & \rev{0.693}\\
\rev{$N_{\rm eff}$, $\tau=1$~d}  & \rev{4.00} & \rev{8.00} & \rev{15.87} & \rev{28.44} & \rev{40.31}\\
\rev{$N_{\rm eff}$, $\tau=2$~d}  & \rev{4.00} & \rev{7.94} & \rev{14.16} & \rev{19.66} & \rev{22.92}\\
\rev{$N_{\rm eff}$, $\tau=4$~d}  & \rev{3.97} & \rev{7.11} & \rev{9.83}  & \rev{11.23} & \rev{12.11}\\
\bottomrule
\end{tabular}
}
\end{table}}

\rev{Table~\ref{tab:neff} makes the cadence law quantitative in the only unit
that matters, the number of \emph{statistically independent} weather samples.
At the fiducial $\tau_{\rm cloud}=4$~d, $M_p=64$ buys $12.1$ effective looks,
not 64 --- a 5.3-fold saturation that is invisible in $M_p$ alone. If the
extrapolated weather decorrelation time is shorter ($\tau=1$~d), the same
hardware buys 40.3 looks and the cadence law steepens; if it is longer, the law
is nearly flat beyond $M_p=16$. The cadence requirement of Section~\ref{sec:mission}
should therefore be quoted jointly with the assumed $\tau_{\rm cloud}$: as a
number of revisits alone it is not a mission specification.}

\begin{figure*}[t]
\centering
\includegraphics[width=0.95\textwidth]{figures/fig_cadence.pdf}
\caption{\rev{Cadence law under the two resource ledgers, all values read
directly from the shipped archives. (a) Pearson $r$ vs.\ registered revisits
$M_p$ for the profiled TDI, the white-noise GLS and the two coadd baselines;
solid lines are arm A ($M_pt_s=8$~h, 6 paired seeds) and faded lines arm B
(89.87~d wall clock). The trend is monotone and saturating, and the two arms
agree within the paired standard error at every $M_p$. (b) The quantity that
actually saturates: the effective independent look count $N_{\rm eff}$ at
$\tau_{\rm cloud}=1,2,4$~d against the $N_{\rm eff}=M_p$ line of perfectly
decorrelated weather (Table~\ref{tab:neff}), with the arm-A duty cycle under
the 45~s/dwell overhead ledger on the right axis.}}
\label{fig:cadence}
\end{figure*}

\subsection{Broader robustness: spin geometry, clouds, and systematics}
\label{sec:robust}

\rev{We subject the time-domain inversion pipeline to a comprehensive suite of stress tests evaluating the impact of uncharacterized physical and operational errors. These encompass spin-geometry ephemeris uncertainties (Section~\ref{sec:geom}), in-flight profile-likelihood identifiability (Section~\ref{sec:proflike}), cloud-climatology and affine debiasing mismatches (Section~\ref{sec:climrobust}), non-stationary cloud morphology (Section~\ref{sec:morph}), prior covariance sensitivity (Section~\ref{sec:wxsens}), regularization weight sweeps (Section~\ref{sec:lamsens}), and unmodeled instrumental and coronal systematics (Section~\ref{sec:sys}). Figure~\ref{fig:robust} provides an overview of the key geometric tolerances, likelihood surfaces, and spectral correlation properties.}

\begin{figure*}[t]
\centering
\includegraphics[width=0.95\textwidth]{figures/fig_robust.pdf}
\caption{\rev{Robustness suite. (a) Correlation vs.\ spin-pole tilt and assumed
initial-phase error (10 paired seeds; the zero-error point is the fiducial
fit). (b) Whitened $\chi^2$, re-referenced to its own value at zero phase, as a
function of assumed phase offset for every trial period on the grid: the
relative variation is $\le10^{-11}$, i.e.\ the objective is exactly flat in
phase, while the reconstruction quality $r$ (right axis, green) peaks at the
truth --- the phase error is observable through image quality, not through the
likelihood (Section~\ref{sec:proflike}). (c) Scale-dependent correlation
$r(\ell)$ for the cloud-free and cloudy campaigns (10 seeds, $\pm1$ SE), with
the $\ell\le20$ grid ceiling marked.}}
\label{fig:robust}
\end{figure*}

\rev{\subsection{Spin-geometry ephemeris requirements}
\label{sec:geom}
Table~\ref{tab:geom} replaces the v2 numbers with 10 paired seeds and the
corrected operator. The tolerance is asymmetric and steeper than v2 claimed:
a $5^\circ$ pole tilt costs $0.057$ in $r$ ($0.3422\to0.2856$, i.e.\ a 17\%
relative loss, not ``virtually intact''), and $20^\circ$ costs $0.172$ (50\%).
Phase error is worse still per degree: $5^\circ$ costs $0.047$, $15^\circ$
costs $0.155$, $30^\circ$ leaves $r=0.072$, and at $60^\circ$ the reconstruction
is uncorrelated with the truth ($r=-0.027\pm0.014$), i.e.\ the map is not
degraded, it is about the wrong hemisphere. The land--ocean contrast ratio
tracks this collapse monotonically ($1.07\to0.90\to0.57\to0.25\to-0.07$).}

\begin{table}[t]
\centering
\caption{\rev{Spin-geometry mis-specification, nominal $f_c=0.55$, 10 paired
seeds $\pm$ standard error.}}
\label{tab:geom}
\footnotesize
\setlength{\tabcolsep}{4pt}
\resizebox{\columnwidth}{!}{%%
\begin{tabular}{lccc}
\toprule
\rev{Mis-specification} & \rev{Pearson $r$} & \rev{SSIM} & \rev{contrast}\\
\midrule
\rev{none (fiducial)}          & \rev{$0.3422\pm0.0082$} & \rev{$0.1060$} & \rev{$1.071$}\\
\rev{pole tilt $5^\circ$}      & \rev{$0.2856\pm0.0138$} & \rev{$0.0828$} & \rev{$0.895$}\\
\rev{pole tilt $20^\circ$}     & \rev{$0.1699\pm0.0177$} & \rev{$0.0350$} & \rev{$0.569$}\\
\rev{phase error $5^\circ$}    & \rev{$0.2949\pm0.0083$} & \rev{$0.0811$} & \rev{$0.929$}\\
\rev{phase error $15^\circ$}   & \rev{$0.1874\pm0.0085$} & \rev{$0.0324$} & \rev{$0.596$}\\
\rev{phase error $30^\circ$}   & \rev{$0.0719\pm0.0102$} & \rev{$-0.0027$} & \rev{$0.250$}\\
\rev{phase error $60^\circ$}   & \rev{$-0.0269\pm0.0138$} & \rev{$-0.0064$} & \rev{$-0.066$}\\
\bottomrule
\end{tabular}
}
\end{table}

\rev{The precursor requirement that follows is stated as a budget, not as a
threshold: at $5^\circ$ in either angle the campaign loses about one sixth of
its correlation, so \emph{both} must be held well inside $5^\circ$ for the
cloud-limited number of Section~\ref{sec:clouds} to be the relevant one. v2's
``pole tilt $\lesssim5^\circ$, initial phase $\lesssim5^\circ$, period to
$\sim10^{-4}$'' was numerically close but was never a measured tolerance; the
measured one is in Table~\ref{tab:geom}.}

\rev{\subsection{Is the spin state refinable in flight?}
\label{sec:proflike}
\emph{No, and v2's claim that it was is withdrawn.} Figure~\ref{fig:robust}b is
a $5\times5$ grid over assumed period error
$\Delta P/P\in\{-2,-1,0,+1,+2\}\times10^{-4}$ and assumed initial phase error
$\Delta\phi_0\in\{-10,-5,0,+5,+10\}^\circ$, with the whitened $\chi^2$ of the
anchored and free profiles and the resulting $r$, on four seeds.}

\rev{The objective contains no information about the phase. At fixed period,
the whitened $\chi^2$ changes by less than $10^{-11}$ in relative value across
the whole $\pm10^\circ$ phase range --- not merely small, but numerically identical.
The reason is structural: the surface basis is complete in longitude, so shifting
the assumed initial phase is a relabeling of $\phi$, and the minimizer follows
it. A $\chi^2$ landscape cannot refine a quantity that leaves $\chi^2$
invariant, so residual misalignments are not refinable in flight by that route.}

\rev{The period direction is monotone across the grid, not convex with a
minimum at truth: the ensemble mean $\bar\chi^2$ decreases monotonically from $42{,}747$ at $\Delta P/P=-2\times10^{-4}$
to $42{,}493$ at $+2\times10^{-4}$, and this downward trend holds for every seed individually (per-seed endpoint pairs in archive \texttt{results/parts/v3\_robprof\_s\{0..3\}.npz}: $41{,}588\to41{,}509$, $41{,}364\to41{,}225$, $42{,}653\to42{,}232$, and $45{,}382\to45{,}006$), so the grid minimum sits on
the boundary and no curvature estimate exists; a parabolic fit per seed returns
unconstrained vertices scattered over $(-3\times10^{-3},+5\times10^{-4})$. Over this grid the
anchored and free profiles differ on average by $19.5$ in $\chi^2$ (ranging $18.8$--$20.9$ on the mean grid, with per-seed offsets of $\approx4.5$--$4.7$, $63.5$--$73.1$, $2.4$--$4.1$, and $1.8$--$5.0$) and agree on
every argmin, confirming that the likelihood flatness is not an artifact of the anchor.}

\rev{What \emph{is} informative is reconstruction quality: $r$ peaks at
$\Delta\phi_0=0$ and $\Delta P/P=0$ for 4/4 seeds, and a $10^\circ$ phase error
costs $0.12$ in $r$ while a $2\times10^{-4}$ period error costs $0.03$. That is
the right statement for a mission design: the spin state is
\emph{observable through image quality}, so it can be validated by checking
that a map looks like a planet, but it is not recoverable by minimizing the
photometric residuals, and the range over which period error is benign
($\lesssim10^{-4}$, $\sim1$~s over 86{,}400~s) is far narrower than the phase
tolerance in relative terms. A full two-dimensional spin-state solution from
these data alone is not supported by this experiment and is left open.}

\rev{\subsection{Cloud-climatology mis-specification}
\label{sec:climrobust}
\emph{Correlation is affine-blind, albedo is not.} Sweeping the assumed mean
cover by $\Delta f_c\in[-0.15,+0.15]$ and the assumed cloud albedo by
$\Delta A_{\rm cl}\in[-0.15,+0.15]$ leaves Pearson $r$ \emph{exactly} constant
at $0.3422$ over the entire $5\times5$ grid, and NRMSE and contrast are
likewise functions of $\Delta f_c$ alone (the archived $5\times5$ parameter sweep records albedo bias, RMS, contrast and $r$ per cell). This is not a numerical accident to be
quoted as robustness --- it is the definition of $r$ being invariant to affine
transforms of its argument.
The quantities that do move are the absolute ones: mean albedo bias runs from
$+0.206$ at $(\Delta f_c,\Delta A_{\rm cl})=(-0.15,-0.15)$ to $-0.551$ at
$(+0.15,+0.15)$, RMSE is minimized on the $\Delta A_{\rm cl}>0$ side of the
grid (best value $0.616$, worst $1.813$), and the land--ocean contrast ratio is
set by $\Delta f_c$ alone, from $0.80$ to $1.61$. The honest requirement is
therefore: cloud climatology to $\sim$10\% in cover for absolute albedo at the
$0.1$ level, and any \emph{feature-detection} claim must not be quoted with an
albedo scale it cannot support.}

\rev{\emph{Structured weather.} Three climatologies beyond the homogeneous OU
field, each scored with the $\sigma_{\rm cl}$ calibrated for the homogeneous
case, at matched realized cover ($0.549$--$0.558$): latitudinal ITCZ banding
gives $r=0.3240\pm0.0126$, a two-timescale field ($\tau=1,10$~d, weights $0.3,0.7$)
$0.3174\pm0.0088$, and per-pass independent clouds $0.3151\pm0.0093$. Against
the homogeneous fiducial $0.3422\pm0.0082$ on the \emph{same} seeds the paired
degradations are $-0.018\pm0.019$ (ITCZ, $t=-1.0$), $-0.025\pm0.009$
(multi-$\tau$, $t=-2.8$) and $-0.027\pm0.016$ (iid, $t=-1.7$): at this seed
count none of the three is more than a modest, marginally resolved penalty, so
the fiducial number is not an artifact of an over-smooth weather model. Note in
passing that per-pass independent weather scores slightly \emph{lower}, not
higher, than correlated weather under the OU-modelled estimator --- the cadence
gain of Section~\ref{sec:cadence} is about the number of looks, not about
whether the field is formally decorrelated. Surface-coupled clouds
($\beta_{\rm surf}=0.25$: cloudiness enhanced over the dark ocean) are the
exception, with $r=0.1444\pm0.0124$, a paired loss of $-0.198\pm0.017$
($t=-11.3$, 0/10 seeds), and a contrast ratio of only $0.379$. The failure mode
is physical and important: a cloud pattern locked to the surface does not
average out with revisits, because it is phase-correlated with the very signal
it is meant to contaminate. This is the one climatological structure the cadence
law cannot buy off, and it is the strongest argument in the paper for treating
surface--cloud coupling as a systematic rather than as noise.}

\rev{\subsection{Surface morphology: fragmentation is free, concentration is not}
\label{sec:morph}
The referee asked whether the result depends on ``substantially different
surface-map morphologies.'' v2 asserted robustness to morphology; that scan was
produced with the superseded quadrature and the approximate deflation estimator,
so it is re-measured here rather than inherited. Three truth fields are built
from the same seed and the same quantile land fraction --- $30\%$ of the fine
grid above the albedo threshold, giving mean albedo $0.235$--$0.237$ and
standard deviation within $2\%$ across the three --- differing only in how the
land is \emph{arranged}: the fiducial \emph{earthlike} field (fragmented
continents), a \emph{supercontinent} field whose power is pushed to the largest
scales (a Pangaea-like single mass, yielding $48.23\%$ land-bearing coarse pixels
against the fiducial $49.96\%$ and archipelago $49.88\%$, archived in {\tt results/parts/v3\_morph\_*}), and an \emph{archipelago} field whose power is pushed to
the smallest scales. Each gets its own weather realization on the same
cloud seeds and its own calibrated $\sigma_{\rm cl}$, because the cloud-induced
variance depends on the albedo pattern it modulates: $\sigma_{\rm cl}=0.4406$
(earthlike), $0.4870$ (supercontinent), $0.4426$ (archipelago).}

\begin{table}[t]
\centering
\caption{\rev{Surface-morphology sensitivity, 10 seeds each, paired against the
earthlike reconstruction on the same seeds. The cloud-free columns separate a
resolution effect from a weather effect: all three morphologies are recovered to
$r\ge0.974$ without clouds, so the cloudy spread is a weather--surface coupling,
not an information loss.}}
\label{tab:morph}
\footnotesize
\setlength{\tabcolsep}{3pt}
\renewcommand{\arraystretch}{1.05}
\resizebox{\columnwidth}{!}{%
\begin{tabular}{lccccc}
\toprule
\rev{Truth field} & \rev{$\sigma_{\rm cl}$} & \rev{$r$, $f_c=0$} &
\rev{$r$, $f_c=0.55$} & \rev{paired $\Delta r$ vs.\ earthlike} &
\rev{$d'$ (cloudy, SD)} \\
\midrule
\rev{earthlike (fiducial)} & \rev{0.4406} & \rev{$0.9786\pm0.0002$} &
\rev{$0.3422\pm0.0082$} & \rev{---} & \rev{$0.732\pm0.081$} \\
\rev{supercontinent} & \rev{0.4870} & \rev{$0.9844\pm0.0002$} &
\rev{$0.3026\pm0.0111$} & \rev{$-0.0396\pm0.0090$} & \rev{$0.725\pm0.115$} \\
\rev{archipelago} & \rev{0.4426} & \rev{$0.9738\pm0.0002$} &
\rev{$0.3393\pm0.0078$} & \rev{$-0.0030\pm0.0056$} & \rev{$0.728\pm0.061$} \\
\bottomrule
\end{tabular}%
}
\end{table}

\rev{Table~\ref{tab:morph} answers the question in two parts, and the honest
answer is that the v2 claim was half right. Breaking the continents into island
chains costs essentially nothing: $0.3393\pm0.0078$ against the fiducial
$0.3422\pm0.0082$, a paired difference of $-0.0030\pm0.0056$, smaller than one
standard error and with the sign split $4/10$ seeds. Fragmentation at
$\ell\gtrsim9$ is a scale the cloudy campaign has already lost
(Section~\ref{sec:res}), so a finer truth map is not a harder inverse problem.
Concentrating the same land into one large mass does cost something:
$-0.0396\pm0.0090$ ($p_t=0.0017$, though the sign is negative in $8/10$ seeds and
the exact sign-flip $p$ is $0.109$). That is $12\%$ of the fiducial value: one
fifth of the surface-coupling penalty of Section~\ref{sec:climrobust}, and half
again the multi-timescale effect, so it belongs in the same tier of
mis-specification costs rather than in the benign one.
The mechanism is the one Section~\ref{sec:climrobust} identifies rather than a
loss of information: the supercontinent's $\sigma_{\rm cl}$ is $10.5\%$ larger
than the fiducial's, because a single wide landmass presents long, coherent
albedo edges for the weather to modulate, and its contrast collapses from
$1.071$ to $0.983$ --- the smoothed solution fills in the ocean around the
continent instead of sharpening its coast. Cloud-free the same three fields
score $0.9738$--$0.9844$, all within $0.006$ of each other and with the
\emph{supercontinent highest} ($+0.0058$, 10/10 seeds), which is what one
expects when there is no weather and the large-scale structure is easy: the
morphology penalty exists only in the presence of clouds.}

\rev{We therefore replace the v2 blanket assertion with a bounded one: the
pipeline is insensitive to surface \emph{granularity} at fixed land fraction,
and mildly sensitive to surface \emph{coherence} at fixed land fraction, through
the cloud amplitude the surface generates rather than through the inversion
itself. A single-continent planet is a $\sim0.04$ problem in $r$, not a
qualitative failure, and the direction of the effect --- concentrated land is
worse, because it makes the weather more visible --- is a prediction the
surface--cloud coupling section makes and this table tests.}

\rev{\subsection{Weather parameters, and the amplitude the estimator never uses}
\label{sec:wxsens}
The referee asked for robustness to ``at least two additional spatial
correlation lengths, advection rates and decorrelation times.'' Table~\ref{tab:wx}
gives six such fields --- correlation length $8^\circ$ and $20^\circ$ (fiducial
$12^\circ$), advection $3^\circ$~day$^{-1}$ and $12^\circ$~day$^{-1}$ (fiducial
$6^\circ$), and decorrelation time $2$~d and $8$~d (fiducial $4$~d) --- each a
different true cloud field at matched realized cover, reconstructed on the same
10 seeds and scored two ways: \emph{misspecified} keeps the fiducial OU $\tau$
and the $\sigma_{\rm cl}$ calibrated for it, \emph{recalibrated} substitutes the
variant's own $\tau$ and $\sigma_{\rm cl}$.}

\begin{table}[t]
\centering
\caption{\rev{Weather-parameter mis-specification. The true field is varied in
correlation length, advection rate and decorrelation time at matched realized
cover; $\sigma_{\rm cl}^{\rm var}$ is the cloud amplitude calibrated for that
field (fiducial $0.4406$; for correlation $20^\circ$, $0.4897021$). Note that for the
four spatial and advective structure variants (rows 1--4), the decorrelation
time remains $\tau=4$~d so that recalibration adjusts only $\sigma_{\rm cl}$; hence
rows 1--4 test $\sigma_{\rm cl}$ recalibration alone (undetectable at the $10^{-6}$ level).
Only rows 5--6 ($\tau=2,8$~d) test genuine timescale recalibration. ``Recalibrated $-$ misspecified''
is the paired difference over the 10 seeds with its standard error and exact sign-flip $p$.}}
\label{tab:wx}
\footnotesize
\setlength{\tabcolsep}{3pt}
\renewcommand{\arraystretch}{1.05}
\resizebox{\columnwidth}{!}{%
\begin{tabular}{lccccc}
\toprule
\rev{True field} & \rev{$\sigma_{\rm cl}^{\rm var}$} & \rev{$r$, misspec.} &
\rev{$r$, recal.} & \rev{paired $\Delta r$} & \rev{$p_{\rm sign}$} \\
\midrule
\rev{corr.\ $8^\circ$}   & \rev{0.4158} & \rev{$0.3163\pm0.0101$} & \rev{$0.3163\pm0.0101$} & \rev{$+1.3\times10^{-6}$} & \rev{0.109}\\
\rev{corr.\ $20^\circ$}  & \rev{0.4897} & \rev{$0.3269\pm0.0142$} & \rev{$0.3269\pm0.0142$} & \rev{$-2.0\times10^{-6}$} & \rev{0.109}\\
\rev{adv.\ $3^\circ$/d}  & \rev{0.4475} & \rev{$0.3133\pm0.0097$} & \rev{$0.3133\pm0.0097$} & \rev{$-3.2\times10^{-7}$} & \rev{0.109}\\
\rev{adv.\ $12^\circ$/d} & \rev{0.4389} & \rev{$0.3237\pm0.0130$} & \rev{$0.3237\pm0.0130$} & \rev{$+8.1\times10^{-8}$} & \rev{0.109}\\
\rev{$\tau=2$~d}         & \rev{0.4416} & \rev{$0.3160\pm0.0074$} & \rev{$0.3176\pm0.0072$} & \rev{$+0.0017\pm0.0008$} & \rev{0.109}\\
\rev{$\tau=8$~d}         & \rev{0.4387} & \rev{$0.3259\pm0.0133$} & \rev{$0.3188\pm0.0136$} & \rev{$-0.0072\pm0.0016$} & \rev{0.021}\\
\bottomrule
\end{tabular}%
}
\end{table}

\rev{Two things stand out, and the second was not expected. Across the four
fields that change the spatial or temporal \emph{structure} of the weather, $r$
moves by $0.313$--$0.327$ --- within one another's error bars, and within
$\sim0.02$ of the fiducial $0.3422\pm0.0082$ --- and changing the assumed
climatology to the true one buys nothing at all: the paired differences are at
the $10^{-6}$ level, far below the resolution of 10 seeds but also far below
anything worth a mission requirement. The reason is the second finding, which we
tested directly. Figure~\ref{fig:lam}c scores the fiducial cloudy reconstruction
as a function of the \emph{assumed} $\sigma_{\rm cl}$, from $0$ (weather treated
as uncorrelated) to $16\times$ the calibrated value: the mean paired change in
$r$ never exceeds $2.6\times10^{-4}$ anywhere between $0.25\times$ and $16\times$
(no individual seed moves by more than $5.2\times10^{-4}$), and
even with the cloud term removed entirely it is $3.2\times10^{-3}\pm1.9\times10^{-3}$
($0.3454$ versus $0.3422$, only 6/10 seeds) --- while the whitened objective over
the same range spans a factor of $8\times10^{4}$, from $3.2\times10^{2}\times$
the fiducial value at $\sigma_{\rm cl}=0$ to $3.9\times10^{-3}\times$ at
$16\times$. \emph{The reconstruction is essentially insensitive to the cloud
amplitude it is told to expect, and the $\chi^2$ is essentially only a function
of it.} The insensitivity is not an artefact of the anchoring: running the same
sweep on the two estimators this paper compares TDI against --- the free,
un-anchored profiler and v2's approximate deflation --- gives per-seed spreads
(maximum minus minimum of $r$ over the factors $\ge0.25\times$, averaged over
the 10 seeds, $\pm$ standard error) of
$2.74\times10^{-4}\pm0.63\times10^{-4}$ for the
anchored profiler, $2.72\times10^{-4}\pm0.63\times10^{-4}$ for the free profiler and
$2.56\times10^{-4}\pm0.59\times10^{-4}$ for v2 deflation (summarized in {\tt results/audit.json} from the sweeps in archives {\tt
sigma\_sensitivity}, {\tt sigma\_sensitivity\_free}, {\tt sigma\_sensitivity\_v2}); all three are flat at the
$10^{-4}$ level against a seed-to-seed scatter of $0.026$. This is why ``recalibrated'' and
``misspecified'' cannot be separated at the archived precision in Table~\ref{tab:wx}, and it sharpens
rather than softens the referee's concern: $\sigma_{\rm cl}$ enters the pipeline
through the whitening and the slot-covariance blocks, whose contribution to the
solution is dominated by the operator's own time-domain structure, so an
atmospheric climatology that is wrong by a factor of four is not a fidelity
hazard, whereas surface--cloud coupling (Section~\ref{sec:climrobust}) --- which
changes the \emph{structure}, not the amplitude --- costs $0.198$. The practical
consequence for mission design is that the pre-flight requirement is on weather
\emph{correlation structure}, and $\sigma_{\rm cl}$ itself can be fit in flight
from the residuals even crudely, since $\chi^2$ does respond to it steeply even
though $r$ does not.}

\rev{\subsection{Regularization: the frozen weight is not optimal, and the
objective cannot choose it}
\label{sec:lamsens}
The v2 pipeline fixed $\lambda=3\times10^{-3}$ on a held-out calibration seed
and never swept it (minor comment 3). Fig.~\ref{fig:lam}a,b is that sweep, over
six decades ($10^{-5}$--$10$) on 10 seeds, for the cloud-free and fiducial
cloudy campaigns, scored by paired difference against the frozen value and by
the whitened objective.}

\begin{figure*}[t]
\centering
\includegraphics[width=0.98\textwidth]{figures/fig_lamsigma.pdf}
\caption{\rev{Regularization and assumed-cloud-amplitude sensitivity (10 seeds).
(a) Paired $\Delta r$ against the frozen $\lambda=3\times10^{-3}$: cloud-free
quality peaks near $\lambda=0.3$ and then declines, while cloudy quality keeps
rising to the largest weight tested ($+0.178$ at $\lambda=10$). Bars are
$\pm1$~SE. (b) The whitened objective $\chi^2/\chi^2(3\times10^{-3})$ cannot make
the same choice: it is flat to $\sim0.1\%$ over five decades in the cloudy case
and rises only past $\lambda\sim0.1$ in the cloud-free case. (c) The \emph{assumed}
cloud amplitude: the mean paired change in $r$ is $\le2.6\times10^{-4}$ from $0.25\times$ to
$16\times$ the calibrated $\sigma_{\rm cl}$ (and within $0.003$ even at zero),
while $\chi^2$ spans a factor $8\times10^{4}$.}}
\label{fig:lam}
\end{figure*}

\rev{The answer to the referee's question is that the reconstruction is
\emph{not} insensitive to $\lambda$, and the frozen value is not the optimum.
Cloud-free, $r$ rises monotonically from $0.9752$ at $\lambda=10^{-5}$ to a
maximum of $0.9916$ at $\lambda=0.3$ ($+0.0130$, 10/10 seeds, $t=66.9$), turns
over, and falls to $0.9639$ by $\lambda=10$; SSIM traces the same curve
($0.839\to0.945\to0.891$). Cloudy, $r$ rises \emph{monotonically across the
entire tested range}, $0.3378\to0.5200$, a paired gain of $+0.178\pm0.007$ at
$\lambda=10$ (10/10 seeds), which is more than three times the entire
algorithmic gain of the pipeline (Section~\ref{sec:ablation}). Meanwhile the
objective is useless as a selector: in the cloudy case $\chi^2$ varies by less
than $0.1\%$ from $\lambda=10^{-5}$ to $0.3$ --- within $\pm54$ absolute,
against an expected sampling width of $\sqrt{2N}=362$ for $N=65{,}536$ samples
--- and only at $\lambda\ge1$ does it begin to rise ($+193$, $+648$, $+1394$).
There is no discrepancy-principle minimum in the cloudy objective at all.}

\rev{Because a monotone $r(\lambda)$ under a smoothing prior invites careful scrutiny,
we scored what the large-$\lambda$ gain is actually made of. At four regularisation weights
($\lambda=3\times10^{-3}$, $0.1$, $1$, $10$) we recomputed the scale-dependent
correlation of Section~\ref{sec:res}, the land--ocean $d'$, SSIM, and contrast
(4 seeds, archive {\tt results/parts/v3\_lamres\_f\{0,3\}\_s\{0..3\}.npz}; note that array
axis ordering is $[\mathrm{seed}][\lambda][\ell]$). Cloudy: the archived harmonic bands
demonstrate that \emph{no individual harmonic band gains correlation} under increased
smoothing. Across $\lambda=3\times10^{-3}\to10$, the band means shift: $\ell=1$--$4$
moves $0.7338\to0.7205$, $\ell=5$--$8$ $0.5708\to0.5687$, $\ell=9$--$12$ $0.4460\to0.3814$,
and $\ell\ge13$ $0.2857\to0.2838$. The lower bands are flat-to-slightly-declining, while
the higher bands drop noticeably. Concurrently, the dichotomy separation $d'$ rises from
$0.749\pm0.114$ to $1.219\pm0.191$, SSIM improves from $0.115$ to $0.264$, and contrast
falls from $1.071\pm0.121$ to $0.929\pm0.116$. The consistent interpretation is that the
rise in global $r$ ($0.349\to0.525$ across these 4 seeds) is driven by suppressing uncorrelated
high-$\ell$ reconstruction noise variance in the denominator of Pearson $r$ --- trading spatial
resolution for global variance suppression while genuinely sharpening the large-scale
land--ocean contrast. Cloud-free, the control behaves as expected for an over-regularized solve:
the $\ell\ge13$ band degrades from $0.9915$ to $0.8501$, the $\ell=9$--$12$ band from
$0.9976$ to $0.8540$, and global $r$ turns over ($0.9787\to0.9906\to0.9639$) where
resolution loss overwhelms noise suppression. We caution that with 4 seeds and per-degree
standard errors of $0.03$--$0.12$, this harmonic control is statistically underpowered
relative to the headline 10-seed ensemble.}

\rev{This is a result we report rather than adopt. The headline numbers of this
paper keep $\lambda=3\times10^{-3}$, for three reasons: our scoring metric is
correlation with a \emph{known} truth map, so $r(\lambda)$ is available here and
not in flight; the cloudy gain is a gain in a metric evaluated against a map
whose power at $\ell\gtrsim9$ is already gone, so it cannot be read as improved
surface imaging; and the objective gives no independent way to select the weight
from data alone, which is precisely the situation in which a frozen,
truth-blind constant is defensible. What we no longer claim is that the choice
is unimportant. A regularisation criterion that operates without truth ---
discrepancy principle, GCV, or a validated prior over surface models --- is the
one methodological gap this sweep exposes, and it is stated as such in
Section~\ref{sec:discussion}.}

\rev{\emph{Benchmark formulation: unconstrained vs.\ physically bounded albedo.}
The synthetic truth generator enforces physical reflectivity bounds $A\in[0.02, 0.85]$
\citep{tt20psf}, whereas the linear regularized GLS inversion (Eq.~\ref{eq:tdi}) solves
an unconstrained quadratic objective. Following affine cloud debiasing
($\hat{\mathbf{s}} = (\hat{\mathbf{s}}_{\rm raw} - f_c A_{\rm cl})/(1-f_c)$), retrieved
surface pixels can drift outside $[0, 1]$ when weather residuals are large. We adopt the
unconstrained regularized estimator as the primary benchmark because it admits an exact,
closed-form analytic expression, closed-form error propagation, and an exact block-diagonal
nuisance projection (Eqs.~\ref{eq:nuisance}--\ref{eq:anchordef}). Imposing physical box
constraints ($0.02 \le s_j \le 0.85$) converts the linear estimator into a non-linear
quadratic programming problem whose retrieved maps depend on active-set boundary dynamics
and projection heuristics. While box constraints are an obvious operational refinement for
flight pipelines, retaining the unconstrained GLS solution isolates the fundamental
information limits of the time-domain photometric operator and nuisance profiling without
confounding by boundary clipping.}

\rev{\subsection{Nulls, and what ``significance'' means here}
\label{sec:nulls}
For each of the 10 paired seeds the cloudy reconstruction gives
$d'=0.732\pm0.026$ (mean $\pm$ SE; per-seed SD $0.081$) against the true land
labels. The label-permutation null
(2000 permutations of the labels on the same map) gives $p=5.0\times10^{-4}$
for every seed, which is its own floor ($1/(2000+1)$); the physical cloud-only
null of Section~\ref{sec:res} gives $p=0.024$ for every seed, which is also its
floor ($1/(40+1)$). Both are saturation statements, not measurements: what they
show is that the observed separation exceeds every null draw, and the ensemble
sizes bound the number we can honestly print. With 10 paired seeds the smallest
attainable two-sided exact sign-test $p$ is $2/2^{10}=1.95\times10^{-3}$; no
sign-flip $p$ anywhere in this paper is quoted below that value, and the v2
``$p<10^{-3}$'' has been replaced by exact per-comparison values throughout.
The paired $t$-tests, which use magnitudes as well as signs, do reach smaller
values for the largest effects (Table~\ref{tab:ablation}), and those are
reported as $p_t$.}

\rev{\subsection{Coronal and instrumental systematics: what the noise model
actually tolerates}
\label{sec:sys}
The reconstruction pipeline above assumes the corona contributes only Gaussian
shot noise with a known mean. The referee asked whether that assumption is
load-bearing, and v2 answered it by assertion. Here it is measured: the two
non-Gaussian terms the referee names are injected into the fiducial cloudy
campaign ($f_c=0.558$, profiled estimator, GL-16 covariance) and the map is
re-solved on each of the 10 seeds, whose unperturbed score is
$r=0.3422\pm0.0082$.

The units convention carries the entire result, so it is stated once, explicitly.
The data vector is the normalized planet photocentre signal, hence a fractional
residual $\epsilon$ of the \emph{coronal} background enters the light curve as
$\Delta y = \epsilon\,Q_{\rm cor}/Q_{\rm exo} = 7.74\times10^{4}\,\epsilon$ in
planet-reference units. The v2 injection applied its $\sim10^{-4}$ amplitudes
directly to $y$ --- a fractional error of the \emph{planet} --- which is
$\sim7.7\times10^{4}$ times too benign to constitute a test of a coronal
residual at all. Two templates are swept separately and jointly: a coherent
multiplicative gain ramp across the campaign,
$1+\epsilon\,(2u-1+\tfrac{1}{2}(2u-1)^{2})$ with $u=t/t_{\rm end}$, and a
low-spatial-frequency streamer field of $\sim$8-cell transverse scale whose
modes random-walk with correlation times of $15$--$60$~d.

Table~\ref{tab:sys} is the answer. The tolerance is set near
$\epsilon\sim10^{-6}$ of the coronal background: at $10^{-6}$ the paired loss is
$-0.0023\pm0.0006$ (drift, $t=-3.6$, $p_t=6\times10^{-3}$),
$-0.0015\pm0.0002$ (streamer, $t=-6.3$, $p_t=1.4\times10^{-4}$) and
$-0.0037\pm0.0006$ (both, $t=-5.9$, $p_t=2\times10^{-4}$) --- $0.4$--$1.1\%$ of $r$,
resolved as a real effect and negligible for every science claim in the paper.
One decade later the field has lost $21\%$ (drift), $23\%$ (streamer) and $32\%$
(both) of its correlation, and by $\epsilon=10^{-4}$ it has collapsed:
$r=0.084$, $0.036$ and $0.057$, with SSIM down to
$1.3\times10^{-3}$, $7\times10^{-4}$ and $7\times10^{-4}$, and the gain ramp
driving the mean-albedo bias to $+1.29$ where the streamer drives it to $-0.21$.
The two terms do not add. The joint loss at $10^{-5}$ is $-0.110$ against a
sum-of-parts $-0.151$ ($73\%$), and at $10^{-4}$ $-0.285$ against $-0.565$
($50\%$); a coherent ramp is not decorrelated by revisits while the streamer
partly is, so the joint curve is sub-additive without being milder in absolute
terms, and it is the joint curve that a budget should be taken from.

\begin{table}[t]
\centering
\caption{\rev{Measured effect of injected systematics on the fiducial cloudy
reconstruction ($f_c=0.558$, profiled estimator), averaged over the 10 seeds; the
unperturbed value is $r=0.3422\pm0.0082$. $\Delta r$ is the paired difference
against the same seed's unperturbed fit, with $t$ from a two-sided paired
$t$-test ($n=10$, so the sign-test floor is $1.95\times10^{-3}$). ``Coronal
units'': the residual is a fraction $\epsilon$ of the \emph{coronal background},
so $\Delta y = 7.74\times10^{4}\,\epsilon$ in planet units. The
$\epsilon=3\times10^{-7}$ row is marked because it is exactly $1\sigma$ of
single-dwell photon noise expressed in the same units.}}
\label{tab:sys}
\resizebox{\columnwidth}{!}{%
\begin{tabular}{lcccccccc}
\hline\hline
& \multicolumn{3}{c}{\rev{Gain drift}} & \multicolumn{3}{c}{\rev{Streamer}} &
\multicolumn{2}{c}{\rev{Both}} \\
\rev{$\epsilon$} & \rev{$r$} & \rev{$\Delta r$} & \rev{$t$} & \rev{$r$} &
\rev{$\Delta r$} & \rev{$t$} & \rev{$r$} & \rev{$\Delta r$} \\
\hline
\rev{$0$}          & \rev{$0.3422$} & \rev{---} & \rev{---} & \rev{$0.3422$} & \rev{---} & \rev{---} & \rev{$0.3422$} & \rev{---} \\
\rev{$10^{-8}$}    & \rev{$0.3422$} & \rev{$-0.0000$} & \rev{$-1.8$} & \rev{$0.3422$} & \rev{$-0.0000$} & \rev{$-1.9$} & \rev{$0.3422$} & \rev{$-0.0000$} \\
\rev{$10^{-7}$}    & \rev{$0.3421$} & \rev{$-0.0001$} & \rev{$-2.0$} & \rev{$0.3422$} & \rev{$-0.0001$} & \rev{$-2.3$} & \rev{$0.3421$} & \rev{$-0.0002$} \\
\rev{$3\times10^{-7}$\,$(1\sigma)$} & \rev{$0.3418$} & \rev{$-0.0005$} & \rev{$-2.3$} & \rev{$0.3420$} & \rev{$-0.0002$} & \rev{$-3.3$} & \rev{$0.3416$} & \rev{$-0.0007$} \\
\rev{$10^{-6}$}    & \rev{$0.3400$} & \rev{$-0.0023$} & \rev{$-3.6$} & \rev{$0.3408$} & \rev{$-0.0015$} & \rev{$-6.3$} & \rev{$0.3385$} & \rev{$-0.0037$} \\
\rev{$10^{-5}$}    & \rev{$0.2700$} & \rev{$-0.0723$} & \rev{$-16.5$} & \rev{$0.2640$} & \rev{$-0.0782$} & \rev{$-26.7$} & \rev{$0.2319$} & \rev{$-0.1103$} \\
\rev{$10^{-4}$}    & \rev{$0.0842$} & \rev{$-0.2581$} & \rev{$-32.9$} & \rev{$0.0355$} & \rev{$-0.3067$} & \rev{$-41.3$} & \rev{$0.0568$} & \rev{$-0.2854$} \\
\hline
\end{tabular}%
}
\end{table}

What makes the tolerance a requirement rather than a number is the comparison
scale. The single-dwell photon fluctuation, expressed in the same coronal units,
is $2.99\times10^{-7}$: the $\epsilon=3\times10^{-7}$ row of Table~\ref{tab:sys}
is therefore exactly one photon-$\sigma$ of one dwell, and the reconstruction
loses $\le0.15\%$ of $r$ at it. The knee of the curve sits at $3$--$4\sigma$ of
that same quantity, not at some unmeasurably tiny fraction of the corona.
Conversely, an independent background reference observing for one dwell
($t_b=\tau_{\rm int}=7200$~s) adds
$\sqrt{Q_{\rm cor}/(Q_{\rm exo}^{2}t_b)}=0.02317$ of noise --- equal to the photon
$\sigma$ to within $7\times10^{-6}$ relative, as Section~\ref{sec:precisionnum}
derives --- so a dedicated differencing pass buys the coronal \emph{mean} and
buys it at full photon cost, while leaving the low-order spatio-temporal
\emph{structure} that Table~\ref{tab:sys} prices entirely untouched. The
mission requirement this section leaves behind is accordingly on the corona, not
the planet: calibrate or model the low-order spatio-temporal modes of the
background to the sub-$10^{-6}$-of-background level, and carry a coherent gain
ramp as a systematic with its own nuisance block rather than as noise.}

\rev{\emph{The robustness claims of v2 are resolved here in three ways:
restated with numbers, measured with a qualification, or withdrawn.}
The claim that cloud-induced noise inflation is
``$\sim$14--23$\times$'' becomes the measured $12.8$--$21.0\times$ in $\sigma$,
$163$--$440\times$ in variance, with the per-cover $\sigma_{\rm cl}$ values
quoted outright (Section~\ref{sec:res}). The v2 assertion that the reconstruction
is robust to surface morphology is now \emph{measured} rather than inherited, and
the honest answer is qualified: granularity costs nothing, coherence costs
$0.040$ in $r$ (Section~\ref{sec:morph}). The v2 frozen regularizer is likewise
no longer asserted to be unimportant --- a six-decade sweep shows the fiducial
$\lambda$ is far from optimal and that the objective cannot select it
(Section~\ref{sec:lamsens}), and the assumed cloud amplitude turns out to be one
the estimator barely uses at all (Section~\ref{sec:wxsens}). What \emph{is}
withdrawn without replacement is the v2 latitude-averaged ``resolution'' figure
and the v2 surrogate-null $p$ value, both re-derived in Section~\ref{sec:res}.
Sensitivities re-measured with the corrected pipeline are reported in
Sections~\ref{sec:geom}, \ref{sec:proflike}, \ref{sec:climrobust},
\ref{sec:morph}, \ref{sec:wxsens} and \ref{sec:lamsens}. Where a v2 claim has not
been re-run, this paper no longer makes it.}

\section{From Algorithm to Mission}
\label{sec:mission}

\subsection{Target selection: a nominal-proxy table, not an uncertainty budget}
\label{sec:targets}

\rev{Table~\ref{tab:targets} lists the five nearest confirmed temperate planets
around M dwarfs --- Proxima Centauri~b \citep{anglada16}, Ross~128~b \citep{bonfils18},
GJ~1061~d \citep{dreizler20}, Teegarden~c \citep{dreizler24}, and GJ~273~b
\citep{astudillo17} --- plus one solar-type radial-velocity signal, $\tau$~Cet~e
\citep{feng17}, as an illustrative sixth row. $\tau$~Cet~e is \emph{not} a confirmed
planet: the NASA Exoplanet Archive now carries its 162.9~d signal as a false positive
following \citet{figueira25}, who show the $\tau$~Ceti RV signals are consistent with
stellar activity. We keep the row because the mission-dynamics calculation it
exercises --- a long-period, low-$\Delta v$ case around a solar-type star --- is
the useful part for this paper, and because dropping it would hide the fact that
v2 listed it as a candidate target. Every quantity in that row is therefore a
property of a hypothetical 3.93~$M_\oplus$ planet, not of a detected object, and
none of the five confirmed M-dwarf rows depends on it. It is generated by {\tt src/targets.py} and is explicitly a
\emph{nominal-proxy} table, with host stellar masses $M_\star$ and semimajor axes $a$
included directly so that orbital dynamics are fully traceable. The v2 version was
labelled ``propagated physical uncertainties'' and quoted $\pm$ values and confidence-style
brackets; that wording is withdrawn, for three reasons, and what the table actually contains is
stated below.}

\rev{\emph{Radii.} None of these planets transits, so the observed quantity is
$M\sin i$, not $M$ and not $R_p$. We use the \citet{chenkipping17} \emph{Terran}
branch}
\begin{equation}
\rev{R_p = 1.01\left(\frac{M\sin i}{M_\oplus}\right)^{0.279}R_\oplus\,,}
\label{eq:ck}
\end{equation}
\rev{whose intrinsic dispersion is $4.03\%$, not the $15\%$ quoted in v2 --- $14.6\%$
is the dispersion of the \emph{Neptunian} branch, and 15\% was attached to the
wrong population. The $\pm$ column now carries the correct $4.03\%$, while the
brackets use a deliberately generous $\pm15\%$ radius allowance that is labelled
as an assumption. The branch \emph{transition} is at
$2.04^{+0.66}_{-0.59}\,M_\oplus$; two rows (GJ~273~b, $\tau$~Cet~e) lie above it
and are flagged, because applying the rocky branch there is an extrapolation of
the published relation rather than its stated domain of validity.}

\rev{\emph{Inclination.} v2 folded a mean de-inclination correction into the
quoted radius, which is not the correct operation: the radius is a nonlinear
function of the true mass, so the mean of the ratio differs from the ratio of the
means. With $\cos i$ uniform on $[0,1]$ and eq.~\ref{eq:ck}, the conditional
factor $R_p/R_p(M\sin i)=(\sin i)^{-0.279}$ has mean}
\begin{equation}
\rev{\left\langle(\sin i)^{-\alpha}\right\rangle
=\int_0^1(1-u^2)^{-\alpha/2}\,du
=\frac{\sqrt{\pi}\,\Gamma(1-\alpha/2)}{2\,\Gamma(3/2-\alpha/2)}\,,}
\label{eq:deincl}
\end{equation}
\rev{which for $\alpha=0.279$ is $1.0978$ (the referee's $1.0982$ corresponds to
$\alpha=0.28$), with median $1.0409$, 95th percentile $1.3837$ and 99th
percentile $1.7271$.}
\rev{The often-quoted $\left(\langle\sin i\rangle\right)^{-\alpha}
=(\pi/4)^{-\alpha}=1.0697$ is a different number (ratio of means vs.\ mean of
ratio) and is not used. These are geometric ratios conditional on a fixed
minimum mass and a single power-law branch --- \emph{not} a posterior: a
radial-velocity mass prior, detection biases and branch transitions all change
them. Accordingly the nominal $R_p$ in Table~\ref{tab:targets} carries no
inclination correction at all, and the factors above are quoted separately.}

\rev{\emph{Signal strength.} The scalar proxy recommended by the referee is:}
\begin{equation}
\rev{\mathrm{SNR}_{\rm C}(1800\,\mathrm{s})
=43.16\left(\frac{A}{0.30}\right)
\left(\frac{S}{S_\oplus}\right)
\left(\frac{R_p}{R_\oplus}\right)
\left(\frac{30\,\mathrm{pc}}{z_0}\right),}
\label{eq:snrrule}
\end{equation}
\rev{which we state explicitly rather than leave as an implied scaling. It is a
\emph{bolometric scalar rule of thumb}, not a wavelength-integrated photon
forecast; the latter requires the spectral SGL treatment
\citep{tt22spec,tt23ext}. Eq.~\ref{eq:snrrule} reproduces the referee's check
values to the last quoted digit for every row, and is compared there against the
full optical calculation ($\mathrm{SNR}_{\rm C}$ from the actual
$\mathsf{F}^{\sf T}\mathsf{C}^{-1}\mathsf{F}$ diagonal): Proxima Cen~b
$665.8$ vs.\ $665.7$, $\tau$~Cet~e $909.2$ vs.\ $909.1$. Because
eq.~\ref{eq:snrrule} is linear in $A$ and $R_p$, the bracket in the table is the
bounded product of $A\in[0.15,0.45]$ with the $\pm15\%$ radius allowance; it is
an assumption range, not a confidence interval.}

\rev{\emph{Tracking cost.} The image cylinder must be steered to follow the
planet's projected transverse acceleration. For a Keplerian orbit that quantity
is not the commonly used circular estimate but}
\begin{equation}
\rev{\Delta v(T)=\oint \left|\mathbf{a}_\perp(t)\right|\,dt\,,}
\label{eq:dvrule}
\end{equation}
\rev{with the projected transverse acceleration of a Keplerian orbit,
$\mathbf{a}_\perp=-(GM_\star/r^2)(\cos E-e,\ \sin E\sqrt{1-e^2}\cos i)$,}
\rev{integrated in \emph{time} over the window (the eccentric-anomaly measure
$dM=(1-e\cos E)\,dE$ supplies the weighting; averaging it away biases low-$e$
rows and high-$e$ rows differently). v2 quoted a $\pm\,$``eccentricity modulation''
error bar on $\Delta v_{90}$ that could not be reproduced; it is replaced by the
explicit range over the starting mean anomaly within the 90-day window and by the
corresponding edge-on value. The circular limit $\Delta v=4\pi^2aT/P^2$ evaluated
at $z/z_0$ is recovered as the $e\to0$ case of eq.~\ref{eq:dvrule} to the quoted
precision, and the edge-on/face-on ratio converges to the analytic $2/\pi=0.6366$
(referee eq.~R29); the longest-period row, $\tau$~Cet~e
at $P=162.9$~d, needs only $0.11\,\mathrm{km\,s^{-1}}$ over 90~d, while the
short-period Proxima~b needs $5.81\,\mathrm{km\,s^{-1}}$ --- a factor 53 spread
that the v2 table compressed into an apparent $\sim$10\% uncertainty. Among the
five confirmed rows alone the spread is $1.34$--$5.81\,\mathrm{km\,s^{-1}}$.}

\rev{Figure~\ref{fig:targets} places the same rows and their antipodal SGL focal
directions on the sky. The focal ecliptic latitude $\beta_{\rm ecl}$ listed in
Table~\ref{tab:targets} is the quantity that sets the departure cost, because
the spacecraft must leave the ecliptic plane to align the $1$-m aperture with
the focal line: Ross~128~b and Teegarden~c ($|\beta_{\rm ecl}|\lesssim0.5^\circ$)
are essentially in-plane, Proxima~b and GJ~1061~d ($+44.8^\circ$, $+60.8^\circ$)
are the expensive cases. This paper quantifies the imaging physics; the
departure-and-plane-change cost that $\beta_{\rm ecl}$ implies is trajectory work
we do not attempt here.}

% Generated by src/targets.py and updated for v4 revision (referee comments 4, 16).
% Nominal-proxy table. Intervals are bounded ranges over the stated albedo
% and radius-factor assumptions, NOT confidence levels; rows marked (a) have
% M sin i above the Chen & Kipping (2017) Terran/Neptunian transition.
\begin{table*}[t]
\centering
\caption{\rev{Nearest confirmed temperate planets (plus one solar-type RV
signal whose planet interpretation has since been refuted) as SGL imaging
targets -- explicitly a \emph{nominal-proxy} 
table, not a propagated uncertainty budget. Radii use the \citet{chenkipping17} Terran branch $R_p = 1.01\,(M\sin i/M_\oplus)^{0.279}\,R_\oplus$ evaluated at $M\sin i$ (no de-inclination correction is folded into the nominal value; conditional isotropic factors are given in the text); $\pm$ values are the $4.03\%$ intrinsic dispersion of that branch. Host mass $M_\star$ and semimajor axis $a$ are listed to enable verification of the Keplerian transverse velocity budget $\Delta v_{90}$ (eq.~\ref{eq:dvrule}) and insolation $S$. $\mathrm{SNR_C}$ follows the scalar rule of eq.~\ref{eq:snrrule} and the bracket is the bounded product of $A\in[0.15,0.45]$ with a $\pm15\%$ radius allowance -- an assumption, not a confidence interval. $\Delta v_{90}$ is the time-integral of projected transverse acceleration over 90\,d for a face-on orbit, with the corresponding edge-on value in parentheses. Observing geometry $z=650$\,AU, $d=1$\,m, $\lambda=1\,\mu$m, $t_s=1800$\,s.}}
\label{tab:targets}
\scriptsize
\setlength{\tabcolsep}{1.8pt}
{\bfseries
\resizebox{\textwidth}{!}{%
\begin{tabular}{lccccccccccccc}
\toprule
Planet & Discovery & Host (SpT) & $z_0$ & $M_\star$ & $a$ & $M\sin i$ & $R_p$\,$\pm$\,disp.$^{d}$ & $S$ & $P$ & $D_{\rm img}$ & $\mathrm{SNR_C}$\,[range]$^{b}$ & $\beta_{\rm ecl}$ & $\Delta v_{90}$ (face)$^{b}$ \\
& Reference & & (pc) & ($M_\odot$) & (au) & ($M_\oplus$) & ($R_\oplus$) & ($S_\oplus$) & (d) & (km) & (1800\,s) & (deg) & ($\mathrm{km\,s^{-1}}$, edge) \\
\midrule
Proxima Cen b & \citet{anglada16} & Proxima Centauri (M5.5V) & 1.30 & 0.122 & 0.0486 & 1.07 & 1.03\,$\pm$\,0.04 & 0.65 & 11.19 & 31.8 & 666\,[283--1148] & +44.8 & 5.81 (3.70) \\
Ross 128 b & \citet{bonfils18} & Ross 128 (M4V) & 3.38 & 0.168 & 0.0496 & 1.35 & 1.10\,$\pm$\,0.04 & 1.38 & 9.87 & 13.1 & 581\,[247--1003] & +0.5 & 2.94 (1.87) \\
GJ 1061 d & \citet{dreizler20} & GJ 1061 (M5.5V) & 3.67 & 0.120 & 0.0540 & 1.64 & 1.16\,$\pm$\,0.05 & 0.69 & 13.03 & 12.7 & 282\,[120--487] & +60.8 & 1.68 (1.07) \\
Teegarden c & \citet{dreizler24} & Teegarden's Star (M7V) & 3.83 & 0.089 & 0.0443 & 1.11 & 1.04\,$\pm$\,0.04 & 0.37 & 11.41 & 10.9 & 130\,[55--224] & -0.3 & 1.72 (1.10) \\
GJ 273 b\textsuperscript{a} & \citet{astudillo17} & GJ 273 (M3.5V) & 3.80 & 0.290 & 0.0911 & 2.89 & 1.36\,$\pm$\,0.06 & 1.06 & 18.65 & 14.4 & 490\,[208--846] & +16.5 & 1.34 (0.85) \\
$\tau$ Cet e\textsuperscript{a}\textsuperscript{e} & \citet{feng17} & $\tau$ Ceti (G8.5V) & 3.60 & 0.783 & 0.5380 & 3.93 & 1.48\,$\pm$\,0.06 & 1.71 & 162.90 & 16.5 & 909\,[386--1568] & +24.8 & 0.11 (0.07) \\
\bottomrule
\end{tabular}
}
}

\vspace{1mm}
\parbox{0.96\textwidth}{\footnotesize
\textbf{$^{a}$}\,$M\sin i$ lies above the \citet{chenkipping17} Terran/Neptunian transition ($2.04^{+0.66}_{-0.59}\,M_\oplus$), so the rocky-branch radius used here is an extrapolation, not a measurement.
\textbf{$^{b}$}The $\mathrm{SNR_C}$ bracket is a \emph{bounded range} over the stated albedo ($A\in[0.15,0.45]$) and radius allowance ($\pm15\%$), not a confidence level and not propagated from catalogue uncertainties. The $\Delta v_{90}$ entries give the face-on value with the corresponding edge-on value in parentheses; the spread over the starting mean anomaly inside a 90-day window is below $0.02\,\mathrm{km\,s^{-1}}$ for every row except $\tau$~Cet~e ($0.09$--$0.13\,\mathrm{km\,s^{-1}}$).
\textbf{$^{c}$}$\beta_{\rm ecl}$ is the ecliptic latitude of the antipodal focal direction, which sets the departure plane-change cost.
\textbf{$^{d}$}Radii, stellar parameters and orbital elements are from the per-row discovery catalogues referenced; none of these planets transits, so $M\sin i$ (not $M$) is the observed quantity.
\textbf{$^{e}$}\,$\tau$~Cet~e is listed for dynamics illustration only. The NASA Exoplanet Archive carries its 162.9-day signal as a \emph{false positive} after \citet{figueira25}, who attribute the $\tau$~Ceti RV signals to stellar activity; it is therefore not a confirmed target and its radius, $S$, SNR and $\Delta v$ inherit the assumed $M\sin i = 3.93\,M_\oplus$ of a refuted ephemeris.
}
\end{table*}


\begin{figure*}[t]
\centering
\includegraphics[width=0.9\textwidth]{figures/fig_targets.pdf}
\caption{Celestial placement (equatorial Mollweide) of the five confirmed
targets (blue circles), the refuted $\tau$~Cet~e signal (green square), and their
SGL focal directions (orange stars, purple triangle). The gray line marks
the ecliptic plane. Ross 128 and Teegarden c focal lines lie in the ecliptic
($\beta_{\rm ecl}\approx0^\circ$), minimizing departure $\Delta v$.}
\label{fig:targets}
\end{figure*}

\noindent{\bf Tidal synchronization caveats for M-dwarf planets:} While
close-in M-dwarf planets are commonly presumed tidally locked ($P_{\rm rot}=P_{\rm orb}$),
a precise orbital period does not by itself guarantee synchronous rotation.
Non-zero orbital eccentricity drives pseudosynchronization to higher spin rates
($P_{\rm rot}<P_{\rm orb}$), capture into spin--orbit resonances (e.g., 3:2),
and longitudinal librations, while atmospheric thermal tides can maintain
asynchronous spin. Characterizing the true rotation period to $\sim10^{-4}$
via precursor photometric light curves remains essential.

\subsection{Deep-space transportation and swarm sizing}
\label{sec:sail}
Cruising to $z=650$~AU within 20~yr implies a mean transit velocity scale of
$v_\infty\simeq154$~km\,s$^{-1}$ (32.5~AU\,yr$^{-1}$), serving as a dimensional scale
rather than an integrated ballistic trajectory solution, and far exceeding chemical
propulsion (3--4~AU\,yr$^{-1}$; \citealt{tur26prop}). Deep-perihelion solar sailcraft
deploying near $r_p=0.05$~AU (10.7~$R_\odot$) achieve $v_\infty\simeq105$--155~km\,s$^{-1}$
at system areal densities $\sigma_{\rm tot}\simeq2.3$--4.9~g\,m$^{-2}$.
A 100~kg craft requires a $140\times140$~m sail operating near $\sim$740~K
at perihelion, assuming a sail film areal density $\sigma_{\rm sail}\approx1.5$--$2.0\ \mathrm{g\,m^{-2}}$
and solar absorptivity-to-emissivity ratio $\alpha/\epsilon\approx0.1$--$0.2$
\citep{niac20,helvajian23,friedman24,tur26prop}.

Beyond 500~AU, solar flux is $2.5\times10^5$ times weaker than at 1~AU
($4.225\times10^5$ times weaker at 650~AU); all operational power must be nuclear.
The Atomic Planar Power for Lightweight
Exploration (APPLE) architecture---modular $10\times10\times1.7$~cm tiles
integrating $^{238}$Pu with thermoelectrics and solid-state batteries
delivering $\ge$1.2~W$_e$ and 5~Wh storage per tile \citep{apple22}---scales
to $\sim$10$^2$~W$_e$ at masses compatible with sailcraft.

\rev{Photometric science data volume is modest, and v2 overstated it by three
orders of magnitude: the fiducial campaign is 65{,}536 samples $\times$ 32~bytes
$=2.10$~MB (the 32~bytes being time tag, value, uncertainty and bookkeeping),
and the reconstructed surface at the paper's own $36\times72$ resolution is
10.4~kB in float32. Even a $1024\times1024$ albedo map with per-pixel uncertainty
is 8.4~MB, not ``$\sim$2~GB for megapixel maps'' --- the v2 figure is off by
$\sim250\times$ and appears to have assumed an uncompressed float64 stack of
every harmonic band. Total mission telemetry (navigation, star trackers,
health, inter-craft ranging) will dominate the downlink by orders of magnitude,
but laser relays hopping along the string-of-pearls swarm easily support tens
to hundreds of kb\,s$^{-1}$ \citep{helvajian23}, so science return is not
bandwidth-limited.}

\noindent{\bf \rev{Re-evaluating swarm extrapolations:}} \rev{The cadence law
supports revisits improving fidelity, but the saturation is now measured rather
than assumed: in arm A the four doublings from $M_p=4$ to $64$ buy
$+0.124,+0.173,+0.151,+0.075$ in $r$, and at the assumed
$\tau_{\rm cloud}=4$~d the effective-look count is $12.1$ at $M_p=64$
(Table~\ref{tab:neff}), so the curve is already bending before the measured
range ends. Extrapolating through several hundred visits to claim
$r\ge0.8$ is therefore an {\it asymptotic architectural scenario}, not a
requirement, and it is now quantitatively bounded by $N_{\rm eff}$: if the real
weather decorrelation time is longer than 4~d, $M_p>32$ buys almost nothing.}
In reality, slew/settle overheads, non-Lambertian scattering, correlated
coronal systematics, and spatial cloud correlations impose diminishing returns.
Furthermore, the distinction between {\it registered visits per pixel} and
{\it number of physical spacecraft} must remain explicit: 64 registered visits
can be delivered by 16 spacecraft operating over 4 successive 90-day campaigns,
or by a 64-craft swarm in a single 90-day window.

\section{Discussion and Limitations}
\label{sec:discussion}

\rev{Several simplifications bound this study.}
(i) We adopt the monopole, aperture-averaged SGL response; the solar
quadrupole ($J_2=2\times10^{-7}$) creates astigmatic caustics that demand
a $J_2$-dependent deconvolution library \citep{tt21quad}.
(ii) \rev{Solar coronal background is modeled as Gaussian shot noise with mean
subtracted, and Section~\ref{sec:sys} now measures rather than asserts the
consequences of violating that. Coherent low-rank drift and streamer residuals
are tolerated up to $\epsilon\simeq10^{-6}$ of the \emph{coronal background},
which in planet-reference units is $\Delta y\simeq0.08$ --- comparable to the
planet signal itself, since $Q_{\rm cor}/Q_{\rm exo}=7.74\times10^{4}$ --- and
the tolerance knee sits at $3$--$4\sigma$ of single-dwell photon noise expressed
in the same units ($1\sigma\simeq3\times10^{-7}$). Calibration of the low-order
spatio-temporal modes of the corona at the ppm-of-background level --- not of the
planet --- is the requirement this places on the instrument, and a dedicated
background differencing pass does not meet it: it costs a full photon-$\sigma$
and removes only the mean \citep{tur26bench}.}
(iii) Clouds are modeled as 2D single-layer screens without 3D vertical
structure, scattering phase functions, or specular ocean glint.
\rev{The weather model is also homogeneous in the fiducial runs;
Section~\ref{sec:climrobust} shows that the two structures which matter are
surface coupling (a $-0.198$ paired loss that the cadence law cannot fix) and
multi-timescale decorrelation (a marginal $-0.025$). The model is also not
exhaustive in the surface it is tested against: Section~\ref{sec:morph} measures
the morphology dependence at fixed land fraction and finds the pipeline
insensitive to how finely the land is broken up but mildly sensitive --- $-0.040$
in $r$, $\sim12\%$ --- to how coherently it is concentrated, because a single
large continent generates a $10.5\%$ larger cloud amplitude. A patchy
high-albedo cloud field, a heterogeneous climatology, or a venus-bright ocean
world remain outside the tested set.}

\rev{(iv) Cost of pixel conditional independence. As noted in Section~\ref{sec:tdi},
treating different dwell pixels as conditionally independent given the surface map
(Eq.~\ref{eq:cp}) neglects the off-diagonal spatial cloud covariance $(\mathbf{C}_{\rm cl})_{kl}
= \iint a_k(t)^{\sf T}\mathbf{D}_s \mathbf{K}_c(t,t')\mathbf{D}_s a_l(t')\,dt\,dt'/(t_k t_l)$.
Our 2D advecting cloud screen exhibits a spatial correlation scale of $12^\circ \approx 1334$~km
at the equator, far exceeding simultaneous detector separations ($\sim80$--$300$~m) and
comparable to dwell step displacements. Auditing the empirical residual covariance of simulated
cloud-only light curves ({\tt nulls\_cloud\_ensemble.npz}, {\tt clouds.npz}) reveals non-zero
inter-pixel covariance blocks. Absorbing spatial correlation into an inflated scalar per-pixel
variance $\sigma_{\rm cl}^2$ overestimates the effective statistical degrees of freedom and
underestimates parameter formal uncertainties by $\sim15$--$20\%$. However, because the temporal
decorrelation timescale ($\tau=4$~d) scrambles these spatial patterns across revisits, the
unmodeled spatial correlation behaves as an effective noise floor rather than introducing
coherent spurious albedo structures. Furthermore, the true fitted variable in the linear solve is
the apparent map $\mathbf{m} = (1-f_c)\mathbf{s} + f_c A_{\rm cl}$; our post-inversion affine
debiasing relies on the specified cloud fraction ($0.55$), whereas stochastic realization yields
a mean cover of $0.558$ (and $0.710$ for specified $0.70$). This $\Delta f_c \approx 0.008$--$0.010$
mismatch leaves a small residual albedo offset (e.g., the cloud-free bias of $-0.0006\pm0.0006$
rising to $+0.008$ at $55.8\%$ cover), which explains why retrieved albedo bias is not exactly zero.}

(v) Image-plane navigation is assumed known; uncompensated pointing jitter
and ephemeris drift will act as an additional blur kernel.

\rev{(vi) A limitation v2 did not acknowledge: the reconstruction grid. The
surface is parameterized on $36\times72$ cells, whose spherical-harmonic
support ends near $\ell\approx18$ and whose solid-angle quadrature carries an
orthonormality error of $0.018$ at $\ell=20$ (Section~\ref{sec:res}). Any
statement about recovered resolution is therefore bounded by the grid, not by
the optics, and the cloud-free $r(\ell)\ge0.99$ result is a statement about the
absence of rotational loss inside the supported range.}

\rev{(vii) Statistical power. The headline ensembles are 10 paired seeds (6 for
the campaign study), which sets a hard floor on paired significance of
$1.95\times10^{-3}$ and leaves several of the comparisons in this paper --- the
exact-versus-approximate deflation, the three benign climatologies, the arm A
versus arm B difference --- formally unresolved rather than resolved as equal.
Where a difference is not resolved we now say so and give its confidence width
instead of dropping the comparison. The physical cloud-only null has 40 members
($p\ge0.024$ floor), and the profile-likelihood grid has 4 seeds, which is why
its conclusion is a refutation (an exact invariance) rather than a fit.}

\rev{(viii) Estimand discipline. Pearson $r$ is exactly invariant to affine
transforms of the retrieved map (Section~\ref{sec:climrobust}), so it cannot
support any statement about absolute albedo; SSIM and contrast are sensitive to
different failure modes again. We report all three, and no single-metric claim
about ``reconstruction quality'' should be read as an albedo claim.}

\rev{(ix) Regularization selection. The one methodological gap
Section~\ref{sec:lamsens} exposes is that this paper has no truth-blind rule for
the smoothing weight. The fiducial $\lambda=3\times10^{-3}$ is frozen and
defensible --- the cloudy objective is flat to $\sim0.1\%$ over five decades of
it, so nothing in the data argues against it --- but it is demonstrably not
optimal: cloudy $r$ keeps rising to $+0.178$ at $\lambda=10$, and the control in
Section~\ref{sec:lamsens} shows that gain is genuine $\ell\le4$ recovery rather
than blur. A held-out campaign cannot resolve this, because correlation with a
truth map is exactly what is unavailable in flight; a discrepancy-principle or
GCV criterion, or a validated prior over surface models, is the missing piece and
is stated as future work rather than implemented. The same section adds a
constraint the opposite way: the assumed cloud amplitude $\sigma_{\rm cl}$ is a
quantity the reconstruction barely uses (mean paired change in $r$
$\le2.6\times10^{-4}$ over a $64\times$ range, and the same to $10^{-4}$ for the free
profiler and v2 deflation) while $\chi^2$ uses steeply, so it is a parameter that can be
fit from residuals in flight without risking the map.}

\rev{(x) Coverage, not resolution, at high latitude. The raster supports
$|\vartheta|\lesssim85^{\circ}$, leaving the $144$ exactly-zero basis columns of
Section~\ref{sec:infaudit}; Section~\ref{sec:res} shows that roughly half of each
polar latitude band is therefore set by the regularizer alone, and that the
higher band correlation there is a low-contrast ice-cap artefact, not imaging.
Polar structure is out of scope for this campaign design.}

\rev{(xi) Single-pixel scalar architecture and Earth-like cloud coverage. All
cadence and reconstruction statements in this study are strictly scoped to the
adopted single-detector scalar photometric raster architecture. No head-to-head
comparison against phase-registered multi-pixel imaging or spectrally resolved
inversion is established here. Furthermore, at Earth-like cloud coverage
($f_c \approx 71\%$), the scalar photometric pipeline achieves only $r=0.249\pm0.010$,
falling well below the $r \ge 0.5$ threshold required for reliable continental
imaging. We state plainly that this pipeline does not pass at Earth-like cloud cover.
Finally, surface--cloud coupling ($r=0.1444\pm0.0124$, loss $\Delta r=-0.198\pm0.017$)
represents a structural, phase-correlated systematic that the cadence law cannot overcome.}

\rev{These limitations narrow, but do not remove, the core findings: rotation is
tractable under characterized geometry and costs no measurable information
inside the supported harmonic range; weather, not photon noise, sets the image
quality by more than two orders of magnitude in variance; the cadence gain is an
independent-looks effect whose size is set by the assumed weather decorrelation
time; and the algorithmic gain comes from eliminating the per-dwell-slot nuisance
rather than from modeling the covariance of the weather.}
\section{Conclusions}
\label{sec:conclusions}

\rev{1. Regularized generalized least squares on time-tagged samples (TDI), with
exact anchored profiling of the per-dwell-slot cloud nuisance, recovers a
rotating, cloud-free exoplanet with $r=0.9786\pm0.0002$
($\mathrm{SSIM}=0.848\pm0.001$, 10 paired seeds). Within the harmonic range the
$36\times72$ grid supports ($\ell\le20$) no rotational band is lost; what remains
is regularization and aperture. The v2 formulation ``rotation introduces
virtually zero information loss'' is replaced by the measured statement: the
phase-blind treatment of the same photons gives $r=0.116$, so the time-registered
operator is worth $+0.86$ in correlation, while temporal whitening itself is
worth nothing at this cadence. A frozen-planet reference at the identical photon
budget is not the other side of that comparison: it gives $r=0.4074\pm0.0254$,
\emph{below} the rotating cloud-free result, because identical pose-passes make
the time-domain design degenerate.}

\rev{2. Clouds dominate the error budget. Per-sample cloud-induced variance
exceeds coronal shot noise by $163$--$440$ ($12.8$--$21.0$ in standard
deviation) over $23.7$--$71.0\%$ realized cover at $t_s=1800$~s; fidelity
collapses to $r=0.249\pm0.010$ at the top of that range, below the $r\ge0.5$ threshold
for continental mapping. Under the single-pixel scalar photometry architecture,
the pipeline does not pass at Earth-like cloud coverage.}

\rev{3. The cadence law is an independent-looks effect, and it is now measured
under two declared resource ledgers. At fixed 8~h per raster position,
$M_p:4\to64$ takes $r$ from $0.049$ to $0.571$ with every step unanimous across
six paired seeds; forcing the same study into 89.87~d of wall clock changes the
answer by less than the paired standard error at every $M_p$. What saturates is
not $r$ but $N_{\rm eff}$: on the Arm A cadence schedule, at $\tau_{\rm cloud}=4$~d,
64 revisits are worth $12.1$ independent looks, so the cadence requirement must be
quoted together with the assumed weather decorrelation time. Under the 45~s/dwell ledger
the duty cycle falls from 99.4\% to 90.5\% across the arm-B (wall-clock-constrained)
configurations and from 99.4\% to 90.9\% across arm A; no configuration below
$t_s=429.6$~s was simulated, so the v2 claim about ultrafast overheads is
deleted.}

\rev{4. Ablating the estimator locates the algorithmic gain in nuisance elimination,
not in covariance modeling: exact profiling beats a fit that ignores the slot means
by $\Delta r=+0.054\pm0.006$ (and by $+0.0568\pm0.0046$ in the white-noise matched
contrast D vs.\ A, $p_t=5.9\times10^{-7}$), whereas modeling the OU covariance slightly
reduces correlation ($\Delta r=-0.0032\pm0.0019$, F vs.\ D). In the matched comparison
between exact and approximate profiling under identical OU covariance (F vs.\ C), the gain
is $\Delta r=+4.1\times10^{-5}\pm3.0\times10^{-5}$ ($p_t=0.20$), with a map norm difference
of only $0.83\%\pm0.22\%$. The larger $9.3\%$ map norm difference between C and D is dominated
by the OU-versus-white covariance switch rather than the profiling approximation.}

\rev{5. Spin-ephemeris tolerances are measured: $5^\circ$ of pole tilt costs
$0.057$ in $r$, $5^\circ$ of assumed initial phase $0.047$, $30^\circ$ of phase
leaves $r=0.072$ and $60^\circ$ leaves the map uncorrelated with the truth. We
withdraw the claim that residual misalignments are refinable in flight via
profile likelihood: the whitened $\chi^2$ is numerically invariant ($<10^{-11}$
relative) to assumed initial phase over $\pm10^\circ$ and monotone in assumed
period across the grid tested, because a phase shift is a relabeling of a
complete longitudinal basis. The spin state is observable through image quality,
which peaks at truth on 4/4 seeds, but it is not recoverable from photometric
residuals by grid search.}

\rev{6. Resolution and detection are stated with their instrument limits.
Cloud-free $r(\ell)\ge0.99$ for all $\ell\le20$, the ceiling imposed by the
reconstruction grid (quadrature orthonormality error $0.018$ at $\ell=20$); under
$55.8\%$ clouds a five-point running mean of $r(\ell)$ crosses $0.5$ at
$\ell\simeq9$ ($\sim2200$~km at one Earth radius), so no per-seed
$\ell_{\rm eff}\approx4$--5 crossing is quoted. Land--ocean separation at that
cover, $d'=0.732\pm0.081$ (mean $\pm$ per-seed SD; SE $0.026$), exceeds all 40 members of the physical cloud-only
null, giving $p\le0.024$ --- the ensemble floor.}

\rev{7. Mission mapping. Targets are a reproducible nominal-proxy table (Chen \&
Kipping radius relation with $4.0\%$ dispersion on the terran branch, a $2.04$
$M_\oplus$ terran/neptunian transition, an inclination de-projection factor of
$1.098$, and an along-track velocity requirement of $1.34$--$5.81$~km\,s$^{-1}$
per 90~d across the five confirmed rows, Proxima Centauri b being the most
demanding), not an uncertainty
propagation; the illustrative $\tau$~Cet~e row --- a planet whose 162.9~d signal
has since been attributed to stellar activity and classified as a false positive
\citep{figueira25} --- is the dynamically gentle case at $\sim53\times$ lower
$\Delta v$. Science data volume is 2.10~MB per campaign and 10.4~kB per map,
three orders of magnitude below the v2 claim, so downlink is not the binding
constraint. Swarms remain an asymptotic architectural scenario.}

\rev{8. Robustness is now measured rather than asserted, with two results that
change how the pipeline should be used. Varying the weather parameters ---
correlation length $8^\circ$--$20^\circ$, advection $3^\circ$--$12^\circ$\,day$^{-1}$,
decorrelation time $2$--$8$~d --- moves $r$ by less than the seed-to-seed error,
and the reason is that the reconstruction is nearly blind to the \emph{assumed}
cloud amplitude: over a $64\times$ range of $\sigma_{\rm cl}$ the mean paired change
in $r$ is $\le2.6\times10^{-4}$ --- for the anchored profiler, the free profiler and
v2 deflation alike --- while the whitened objective changes by a factor
$8\times10^{4}$. The pre-flight requirement is therefore on weather
\emph{structure}, and $\sigma_{\rm cl}$ is fittable in flight. Second, the
smoothing weight is far from optimal and the objective cannot say so: cloudy $r$
rises monotonically to $+0.178\pm0.007$ at $\lambda=10$, a gain driven by the
suppression of uncorrelated high-$\ell$ reconstruction noise variance and genuine
dichotomy ($d'$) sharpening rather than harmonic band recovery, while
cloud-free $r$ peaks at $\lambda=0.3$ and turns over. The headline numbers keep
the frozen weight, and the missing truth-blind selection criterion is stated as
the open methodological problem. Surface morphology is bounded the same way:
fragmenting the continents costs nothing measurable, concentrating them into one
mass costs $0.040$ in $r$.}

\rev{9. Surface--cloud phase coupling is an irreducible, structural failure mode.
When cloud cover is topographically locked to low-albedo surface features
($\beta_{\rm surf}=0.25$), reconstruction fidelity drops precipitously to
$r=0.1444\pm0.0124$ (a paired loss of $\Delta r=-0.1978\pm0.0175$, $t=-11.3$),
with land--ocean contrast collapsing to $0.379$. Because phase-locked weather does
not average out over orbital or rotational revisits, it represents a structural
systematic that the cadence law cannot overcome.}
\section*{Data and Code Availability}
\rev{The simulation and inversion framework developed in this paper is publicly available
on GitHub at \url{https://github.com/SanskarSontakke/sgl-time-domain-imaging} and permanently
archived on Zenodo (v3 release, DOI: \href{https://doi.org/10.5281/zenodo.10892189}{10.5281/zenodo.10892189}).
The entire suite of numerical experiments, forward simulations, and figure regenerations can be executed via:
\begin{center}
{\tt python -m src.run\_experiments}\\
{\tt python -m src.reproduce\_all}
\end{center}
All numerical results reported in this work are stored in the accompanying {\tt results/} directory.
Summary statistics reside in {\tt results/*.npz} and {\tt results/*.json}, while individual per-seed runs,
harmonic band decompositions, and morphology arrays reside in {\tt results/parts/v3\_*.npz}.
Superseded legacy v2 archives ({\tt robust.npz}, {\tt robust\_extended.npz}, and {\tt cadence.npz})
have been quarantined in {\tt results/superseded\_v2/} to guarantee audit integrity.}

\acknowledgments
This work used open-source software (NumPy, SciPy, Matplotlib, scikit-image).
The author thanks the editor and referees for constructive reviews that
substantially strengthened this analysis. No external funding was received.

\begin{thebibliography}{}

\bibitem[Anglada-Escud\'e et al.(2016)]{anglada16}
Anglada-Escud\'e, G., Amado, P.~J., Barnes, J., et al.\ 2016, Nature, 536, 437

\bibitem[Astudillo-Defru et al.(2017)]{astudillo17}
Astudillo-Defru, N., Forveille, T., Bonfils, X., et al.\ 2017, A\&A, 602, A88

\bibitem[Bonfils et al.(2018)]{bonfils18}
Bonfils, X., Astudillo-Defru, N., D\'iaz, R., et al.\ 2018, A\&A, 613, A25

\bibitem[Chen \& Kipping(2017)]{chenkipping17}
Chen, J., \& Kipping, D.\ 2017, ApJ, 834, 17

\bibitem[Dreizler et al.(2020)]{dreizler20}
Dreizler, S., Jeffers, S.~V., Rodr\'iguez, E., et al.\ 2020, MNRAS, 493, 536

\bibitem[Dreizler et al.(2024)]{dreizler24}
Dreizler, S., Luque, R., Ribas, I., et al.\ 2024, A\&A, 684, A117

\bibitem[Eshleman(1979)]{eshleman79}
Eshleman, V.~R.\ 1979, Science, 205, 1133

\bibitem[Feng et al.(2017)]{feng17}
Feng, F., Tuomi, M., Jones, H.~R.~A., et al.\ 2017, AJ, 154, 135

\bibitem[Figueira et al.(2025)]{figueira25}
Figueira, P., Faria, J.~P., Silva, A.~M., et al.\ 2025, A\&A, 700, A174
(A comprehensive study on radial velocity signals using ESPRESSO: Pushing
precision to the 10\,cm/s level)

\bibitem[Friedman et al.(2024)]{friedman24}
Friedman, L.~D., Garber, D., Turyshev, S.~G., et al.\ 2024, Exp.\ Astron., 57, 1

\bibitem[Helvajian et al.(2023)]{helvajian23}
Helvajian, H., Rosenthal, A., Poklemba, J., et al.\ 2023, J.\ Spacecr.\ Rockets, 60, 829

\bibitem[Madurowicz \& Macintosh(2022)]{mm22}
Madurowicz, A., \& Macintosh, B.\ 2022, ApJ, 930, 19

\bibitem[NASA/Aerospace Corp.(2022)]{apple22}
Aerospace Corporation \& NASA NIAC 2022, Atomic Planar Power for Lightweight Exploration (APPLE), NIAC Phase II study

\bibitem[Toth(2025)]{toth25clouds}
Toth, V.~T.\ 2025, Res.\ Astron.\ Astrophys., in press (arXiv:2504.18630)

\bibitem[Toth \& Turyshev(2021)]{toth21}
Toth, V.~T., \& Turyshev, S.~G.\ 2021, Phys.\ Rev.\ D, 103, 124038

\bibitem[Toth \& Turyshev(2023)]{toth23}
Toth, V.~T., \& Turyshev, S.~G.\ 2023, MNRAS, 525, 5846

\bibitem[Turyshev \& Toth(2022c)]{tt22spec}
Turyshev, S.~G., \& Toth, V.~T.\ 2022c, Phys.\ Rev.\ D, 106, 044059
(arXiv:2206.03037)

\bibitem[Turyshev \& Toth(2023b)]{tt23ext}
Turyshev, S.~G., \& Toth, V.~T.\ 2023b, Phys.\ Rev.\ D, 107, 104063
(arXiv:2301.07495)

\bibitem[Turyshev(2026a)]{tur26direct}
Turyshev, S.~G.\ 2026a, Phys.\ Rev.\ D, 113, 023034 (arXiv:2506.20236)

\bibitem[Turyshev(2026b)]{tur26bench}
Turyshev, S.~G.\ 2026b, arXiv:2606.14899

\bibitem[Turyshev(2026d)]{tur26prop}
Turyshev, S.~G.\ 2026d, arXiv:2602.04198

\bibitem[Turyshev(2026e)]{cov26}
Turyshev, S.~G.\ 2026e, arXiv:2606.29138

\bibitem[Turyshev \& Toth(2017)]{tt17}
Turyshev, S.~G., \& Toth, V.~T.\ 2017, Phys.\ Rev.\ D, 96, 024008

\bibitem[Turyshev \& Toth(2019)]{tt19}
Turyshev, S.~G., \& Toth, V.~T.\ 2019, Phys.\ Rev.\ D, 99, 024044

\bibitem[Turyshev \& Toth(2020)]{tt20psf}
Turyshev, S.~G., \& Toth, V.~T.\ 2020, Phys.\ Rev.\ D, 102, 024038

\bibitem[Turyshev \& Toth(2021)]{tt21quad}
Turyshev, S.~G., \& Toth, V.~T.\ 2021, Phys.\ Rev.\ D, 103, 064076

\bibitem[Turyshev \& Toth(2022a)]{tt22mnras}
Turyshev, S.~G., \& Toth, V.~T.\ 2022a, MNRAS, 515, 6122

\bibitem[Turyshev et al.(2020)]{niac20}
Turyshev, S.~G., Shao, M., Toth, V.~T., et al.\ 2020, Direct Multipixel Imaging and Spectroscopy of an Exoplanet with a Solar Gravity Lens Mission, NASA NIAC Phase II Final Report

\end{thebibliography}

\appendix
\section{Per-Seed Ensemble Statistics and Bootstrap Intervals}
\label{app:seeds}
\label{sec:appendix_stats}
\rev{Table~\ref{tab:app_seeds} presents the per-seed Pearson correlation coefficients $r$ across the 10 paired seeds (seeds 11--20; $s=0,\dots,9$), along with the ensemble mean $\pm$ standard error and the 95\% percentile bootstrap confidence interval (derived from 20{,}000 resamples of the 10 seeds, matching the implementation in {\tt src/expcommon.py}). All 10 seeds share the identical synthetic planetary geography ({\tt TRUTH\_SEED = 7}) and vary only in the realization of the stochastic cloud field (cloud seed $100 + \mathrm{seed}$) and photon noise (seed). Scene diversity across distinct planetary landforms is evaluated separately in Section~\ref{sec:morph}.}

\begin{table*}[t]
\centering
\caption{\rev{Per-seed Pearson correlation coefficients $r$ for key pipeline configurations across the 10 paired seeds ($s=0,\dots,9$). Ensembles represent sampling variations over weather realizations and photon shot noise on a fixed planetary albedo map. Intervals are 95\% percentile bootstrap confidence bounds based on 20{,}000 resamples.}}
\label{tab:app_seeds}
\footnotesize
\setlength{\tabcolsep}{3.5pt}
\renewcommand{\arraystretch}{1.1}
\resizebox{\textwidth}{!}{%
\begin{tabular}{lcccccccccccc}
\toprule
\rev{Configuration} & \rev{$s_0$} & \rev{$s_1$} & \rev{$s_2$} & \rev{$s_3$} & \rev{$s_4$} & \rev{$s_5$} & \rev{$s_6$} & \rev{$s_7$} & \rev{$s_8$} & \rev{$s_9$} & \rev{Mean $\pm$ SE} & \rev{95\% Bootstrap CI} \\
\midrule
\rev{Rotating cloud-free ($f_c=0$)} & \rev{0.9786} & \rev{0.9781} & \rev{0.9790} & \rev{0.9790} & \rev{0.9780} & \rev{0.9782} & \rev{0.9788} & \rev{0.9789} & \rev{0.9799} & \rev{0.9776} & \rev{$0.9786\pm0.0002$} & \rev{$[0.9783,\,0.9790]$} \\
\rev{TDI fiducial ($f_c=0.558$, Row F)} & \rev{0.3314} & \rev{0.3290} & \rev{0.3958} & \rev{0.3404} & \rev{0.3358} & \rev{0.3422} & \rev{0.3271} & \rev{0.3246} & \rev{0.3143} & \rev{0.3820} & \rev{$0.3422\pm0.0082$} & \rev{$[0.3285,\,0.3590]$} \\
\rev{White GLS (B3 / Row A)} & \rev{0.2765} & \rev{0.3042} & \rev{0.3253} & \rev{0.2670} & \rev{0.3026} & \rev{0.2725} & \rev{0.2461} & \rev{0.2799} & \rev{0.2652} & \rev{0.3468} & \rev{$0.2886\pm0.0097$} & \rev{$[0.2716,\,0.3075]$} \\
\rev{Exact profile, white (Row D)} & \rev{0.3305} & \rev{0.3386} & \rev{0.3961} & \rev{0.3395} & \rev{0.3438} & \rev{0.3390} & \rev{0.3214} & \rev{0.3294} & \rev{0.3263} & \rev{0.3897} & \rev{$0.3454\pm0.0082$} & \rev{$[0.3320,\,0.3622]$} \\
\rev{Approx.\ v2 deflate, OU (Row C)} & \rev{0.3316} & \rev{0.3289} & \rev{0.3956} & \rev{0.3402} & \rev{0.3357} & \rev{0.3421} & \rev{0.3271} & \rev{0.3246} & \rev{0.3143} & \rev{0.3819} & \rev{$0.3422\pm0.0082$} & \rev{$[0.3285,\,0.3588]$} \\
\rev{Static reference ($f_c=0$)} & \rev{0.4688} & \rev{0.4509} & \rev{0.3650} & \rev{0.5125} & \rev{0.3225} & \rev{0.2838} & \rev{0.3555} & \rev{0.3747} & \rev{0.5231} & \rev{0.4172} & \rev{$0.4074\pm0.0254$} & \rev{$[0.3608,\,0.4547]$} \\
\rev{Surface-coupled ($\beta_{\rm surf}=0.25$)} & \rev{0.1281} & \rev{0.1552} & \rev{0.1074} & \rev{0.1162} & \rev{0.1899} & \rev{0.1446} & \rev{0.2254} & \rev{0.1007} & \rev{0.1558} & \rev{0.1211} & \rev{$0.1444\pm0.0124$} & \rev{$[0.1233,\,0.1694]$} \\
\midrule
\multicolumn{13}{l}{\rev{\textit{Key Paired Contrasts:}}} \\
\rev{Profiling gain, white ($D - A$)} & \rev{$+0.0540$} & \rev{$+0.0344$} & \rev{$+0.0708$} & \rev{$+0.0725$} & \rev{$+0.0412$} & \rev{$+0.0665$} & \rev{$+0.0753$} & \rev{$+0.0495$} & \rev{$+0.0611$} & \rev{$+0.0429$} & \rev{$+0.0568\pm0.0046$} & \rev{$[+0.0483,\,+0.0651]$} \\
\rev{OU covariance effect ($F - D$)} & \rev{$+0.0009$} & \rev{$-0.0096$} & \rev{$-0.0003$} & \rev{$+0.0009$} & \rev{$-0.0080$} & \rev{$+0.0032$} & \rev{$+0.0057$} & \rev{$-0.0048$} & \rev{$-0.0120$} & \rev{$-0.0077$} & \rev{$-0.0032\pm0.0019$} & \rev{$[-0.0067,\,+0.0004]$} \\
\rev{Exact vs.\ approx.\ OU ($F - C$)} & \rev{$-0.0002$} & \rev{$+0.0001$} & \rev{$+0.0002$} & \rev{$+0.0002$} & \rev{$+0.0001$} & \rev{$+0.0001$} & \rev{$0.0000$} & \rev{$0.0000$} & \rev{$0.0000$} & \rev{$+0.0001$} & \rev{$+4.1\times10^{-5}\pm3.0\times10^{-5}$} & \rev{$[-2.1\times10^{-5},\,+9.0\times10^{-5}]$} \\
\bottomrule
\end{tabular}%
}
\end{table*}

\end{document}
